% Created 2026-07-09 qui 23:16
% Intended LaTeX compiler: lualatex
\documentclass[12pt,a4paper]{article}
\usepackage{fgv-paper}
% Bibliografia controlada exclusivamente pelo arquivo BibTeX principal.
\usepapercover
\institution{Fundação Getúlio Vargas}
\programname{}
\coursename{Mestrado Acadêmico em Políticas Públicas e Governo}
\disciplinename{Decisões Baseadas em Evidência}
\professorname{Marcos Aurélio Pereira Valadão}
\cityname{Brasília}
\papertype{Artigo acadêmico}
\covernote{Trabalho acadêmico elaborado para a disciplina Decisões Baseadas em Evidência.}
\setlength{\headheight}{15pt}


\author{Gustavo M. Mendes de Tarso}
\date{\today}
\title{ATESTMED, saúde digital e decisão baseada em evidências: revisão estruturada e proposta de redesenho do fluxo pericial}
\hypersetup{
 pdfauthor={Gustavo M. Mendes de Tarso},
 pdftitle={ATESTMED, saúde digital e decisão baseada em evidências: revisão estruturada e proposta de redesenho do fluxo pericial},
 pdfkeywords={},
 pdfsubject={},
 pdfcreator={Emacs 28.2 (Org mode 9.5.5)}, 
 pdflang={Pt_Br}}
\usepackage[backend=biber,style=abnt,language=brazil,sorting=nyt,giveninits=true,uniquename=false,uniquelist=false]{biblatex}
\addbibresource{artigo_final_atestmed_abnt.bib}
\begin{document}

\makemytitle

\begingroup
\small
\begin{center}
\textbf{Resumo}
\end{center}
Este trabalho examina o uso do ATESTMED na Perícia Médica Federal, integrando saúde digital e decisões baseadas em evidências para a concessão de benefícios por incapacidade. O ATESTMED é tratado como análise documental digital assíncrona, distinta da teleperícia, que pressupõe interação remota síncrona, e da perícia presencial, que permite avaliação direta e exame físico. A partir de uma revisão estruturada de 20 estudos com versão integral recuperada, realizou-se análise qualitativa, com descrição quantitativa do fluxo de seleção, identificando quatro eixos centrais — suficiência documental, mediação digital, capacidade operacional e auditoria. A proposta operacional modela um fluxo decisório escalonado, composto por análise documental estruturada pelo ATESTMED, encaminhamento para teleperícia em casos de incerteza intermediária e escalonamento para perícia presencial em situações complexas, contemplando auditoria contínua e indicadores de desempenho por modalidade. A proposta visa garantir eficiência, segurança, auditabilidade e equidade territorial, atendendo desafios como filas, qualidade decisória e desigualdades regionais. A discussão ressalta os riscos da automação parcial, a importância da governança robusta e a necessidade de formação permanente. Limitações metodológicas e necessidade de validação empírica futura são reconhecidas. Conclui-se que a integração criteriosa entre análise documental, teleperícia e perícia presencial, alinhada a indicadores, pode otimizar o processo decisório na Perícia Médica Federal, promovendo sustentabilidade e justiça no sistema previdenciário.

\noindent\textbf{Palavras-chave:} ATESTMED; perícia médica; documentos médicos; decisão baseada em evidências; saúde digital.

\vspace{0.8em}
\begin{center}
\textbf{Abstract}
\end{center}
This paper examines the use of ATESTMED in the Brazilian Federal Medical Expert Service, integrating digital health and evidence-based decision-making in the granting of temporary incapacity benefits. ATESTMED is treated as asynchronous digital documentary analysis, distinct from remote medical expert assessment, which presupposes synchronous remote interaction, and from in-person medical expert examination, which allows direct assessment and physical examination. Based on a structured review of 20 studies whose original full versions were retrieved, a qualitative analysis was conducted, with a quantitative description of the selection flow, identifying four central axes: documentary sufficiency, digital mediation, operational capacity, and auditing. The operational proposal models a tiered decision-making flow composed of structured documentary analysis through ATESTMED, referral to remote medical expert assessment in cases of intermediate uncertainty, and escalation to in-person medical expert examination in complex situations, including continuous auditing and performance indicators by modality. The proposal seeks to ensure efficiency, safety, auditability, and territorial equity, addressing challenges such as queues, decision quality, and regional inequalities. The discussion highlights the risks of partial automation, the importance of robust governance, and the need for continuous training. Methodological limitations and the need for future empirical validation are acknowledged. The paper concludes that the careful integration of documentary analysis, remote medical expert assessment, and in-person medical expert examination, aligned with indicators, can optimize decision-making in the Brazilian Federal Medical Expert Service, promoting sustainability and justice in the social security system.

\noindent\textbf{Keywords:} ATESTMED; medical expert assessment; medical documents; evidence-based decision-making; digital health.
\endgroup


\section{Introdução}
\label{sec:org6c1f968}


A transformação digital do Estado tem reconfigurado, de modo particularmente intenso, os arranjos institucionais de acesso a direitos sociais, os mecanismos de triagem administrativa e os regimes de produção de prova utilizados na tomada de decisão pública. No campo previdenciário, essa inflexão assume relevância singular porque incide sobre benefícios cuja concessão depende da apreciação técnico-pericial de situações de incapacidade laborativa, isto é, de condições clínicas, funcionais e ocupacionais que não se deixam reduzir, sem perdas substantivas, a uma lógica puramente documental ou automatizada. Nesse contexto, o ATESTMED emerge como dispositivo central da agenda recente de saúde digital aplicada à Previdência Social brasileira, ao propor uma sistemática de análise documental digital assíncrona de atestados e documentos médicos para fins de reconhecimento administrativo de incapacidade temporária. A importância do tema decorre não apenas de sua dimensão operacional, associada ao enfrentamento de filas, tempos de espera e gargalos de atendimento, mas também de seus efeitos sobre a qualidade epistêmica da decisão administrativa, sobre a equidade no acesso ao benefício e sobre a própria legitimidade da Perícia Médica Federal como instância de julgamento técnico em larga escala \parencite{soares20241031t0000000000asistematicadoatestmednaconcessa}.

O problema público que se coloca, portanto, não é simplesmente o de informatizar fluxos ou digitalizar documentos. Trata-se de compreender em que medida a adoção de soluções de saúde digital, quando transpostas para o domínio da perícia previdenciária, altera os critérios de suficiência probatória, os padrões de consistência decisória e os mecanismos de controle de qualidade em contextos de alta demanda. A concessão de benefícios por incapacidade envolve, simultaneamente, proteção social, sustentabilidade fiscal, integridade procedimental e responsividade administrativa. Qualquer inovação que amplie escala e velocidade pode produzir ganhos relevantes de acesso, mas também pode introduzir novos riscos, como assimetrias informacionais, exclusão digital, heterogeneidade interpretativa, incentivos organizacionais distorcivos e enfraquecimento da avaliação contextualizada do caso concreto. A literatura internacional sobre telemedicina e serviços digitais em saúde não é tomada aqui como evidência direta de validade do ATESTMED, pois este opera predominantemente por análise documental assíncrona. Ela é mobilizada de forma indireta e complementar para iluminar os limites da mediação tecnológica em decisões clínicas e administrativas, especialmente quanto à qualidade da informação, à infraestrutura tecnológica, à segurança, à aceitabilidade por usuários e profissionais e à desigualdade de acesso \parencite{ftouni2022challengesoftelemedicineduringthecovid19pandemica}. Embora a perícia médica previdenciária não se confunda com a assistência clínica, ela compartilha com esse universo a dependência de registros confiáveis, de mediações tecnológicas robustas e de critérios transparentes de decisão.

No caso brasileiro, a relevância previdenciária do ATESTMED é ampliada pelo volume estrutural de requerimentos por incapacidade temporária e pela histórica tensão entre capacidade estatal de processamento e demanda social reprimida. Em sistemas massivos, a fila não é apenas um problema gerencial; ela se converte em variável distributiva, pois o tempo de espera afeta renda, tratamento, vínculo laboral e confiança institucional. A introdução de uma via documental digital assíncrona pode, nesse sentido, representar resposta pragmática a um cenário de congestionamento decisório. Contudo, a literatura sobre desempenho organizacional adverte que metas de celeridade e expansão da capacidade de processamento, quando não acompanhadas de salvaguardas substantivas, podem induzir disfunções de desempenho, deslocando a atenção do mérito da decisão para indicadores de produção \parencite{kelman2009performanceimprovementandperformancedysfunctionanemp}. Em matéria previdenciária, esse risco é especialmente sensível, porque a decisão não versa sobre um serviço padronizado, mas sobre o reconhecimento de uma contingência social juridicamente protegida, cuja aferição depende de juízo técnico situado entre medicina, funcionalidade e capacidade para o trabalho.

A literatura sobre incapacidade para o trabalho oferece elementos importantes para situar essa complexidade. Estudos clássicos e contemporâneos mostram que a incapacidade não é um dado exclusivamente biomédico, mas uma construção técnico-institucional que articula diagnóstico, funcionalidade, contexto ocupacional, trajetória laboral, incentivos do sistema e normas administrativas \parencite{dijkstrasdatamakingupincapacityforwork, sainsbury2006routesontoincapacitybenefitsfindingsfromqualitati, corden2001incapacitybenefitsandworkincentives}. Diretrizes voltadas ao manejo de afastamentos prolongados e incapacidade de longa duração também enfatizam a necessidade de abordagens integradas, capazes de considerar fatores clínicos, psicossociais e laborais, em vez de se limitar à documentação médica formal da condição de saúde \parencite{gabbay2011niceguidanceonlongtermsicknessandincapacity}. Essa constatação é decisiva para o debate sobre o ATESTMED, pois, se a incapacidade previdenciária é multidimensional, a substituição, ainda que parcial, do exame pericial presencial por análise documental digital assíncrona exige critérios metodológicos rigorosos para definir em quais situações o documento médico é suficiente, em quais é apenas indiciário e em quais a avaliação direta ou a teleperícia síncrona podem ser necessárias.

Além disso, a literatura sobre comportamento avaliativo de médicos peritos e seguradores demonstra que a decisão em incapacidade é sensível a variações de julgamento, repertórios profissionais e instrumentos de mensuração \parencite{steenbeek2011thedevelopmentofinstrumentstomeasuretheworkdisa}. Isso significa que a digitalização do processo não elimina a subjetividade técnica; ao contrário, pode deslocá-la para novas etapas, como a leitura do atestado, a interpretação da completude documental, a inferência sobre limitação funcional e a definição de necessidade de exame complementar. Em outras palavras, a mediação tecnológica não suprime o problema da discricionariedade pericial, mas o reconfigura. A questão central passa a ser se o novo arranjo institucional produz maior padronização baseada em evidências ou apenas uma nova forma de variabilidade decisória, agora menos visível ao usuário e mais dependente da qualidade dos registros submetidos.

Esse ponto se conecta diretamente ao debate sobre saúde digital e inclusão. A literatura recente evidencia que o uso de telemedicina e de plataformas digitais por pessoas com deficiência ou com necessidades complexas é condicionado por fatores de acesso tecnológico, letramento digital, suporte familiar, renda, conectividade e desenho da interface institucional \parencite{albasheer2026telemedicineutilizationanddigitalinclusionamongpe, m2024astudyofdisabilityseveritybarriersandfacilitatingfactor}. Embora telemedicina e ATESTMED sejam procedimentos distintos, esses achados são úteis por analogia parcial, pois mostram que a mediação digital pode redistribuir barreiras de acesso e de produção da prova. Em contextos de vulnerabilidade, a promessa de desburocratização pode conviver com novas barreiras de entrada, especialmente quando a produção e o envio de documentação adequada dependem de recursos desigualmente distribuídos. A experiência internacional com documentação médico-administrativa de deficiência e acesso a benefícios também sugere que a existência de um documento formal não garante, por si só, fruição efetiva de direitos, seja por entraves administrativos, seja por inadequação entre o documento médico-administrativo e os critérios do sistema de proteção \parencite{rajasulochana2026doeshavingadisabilitycertificateensurebenefit}. Para a Perícia Médica Federal, isso implica reconhecer que a qualidade da decisão baseada em documentos não depende apenas do algoritmo procedimental, mas da ecologia social de produção desses documentos.

A discussão torna-se ainda mais relevante quando se observa que diferentes campos da avaliação em saúde e funcionalidade têm investido em classificações, instrumentos e aplicações digitais para qualificar a mensuração de condições complexas \parencite{feudtner2014pediatriccomplexchronicconditionsclassificationsyst, guo2017assessmentofaphasiaacrosstheinternationalclassification}. Tais iniciativas mostram que a digitalização pode ampliar consistência, comparabilidade e rastreabilidade, desde que ancorada em modelos conceituais robustos e em validação empírica. Entretanto, também revelam que a tradução de fenômenos clínicos e funcionais para formatos estruturados exige escolhas classificatórias que podem simplificar excessivamente a realidade. No domínio previdenciário, essa tensão é particularmente aguda, porque a administração necessita de critérios operacionais escaláveis, enquanto o objeto da decisão — a incapacidade para o trabalho — resiste à padronização absoluta. A literatura sobre reabilitação, retorno ao trabalho e avaliação de capacidade laboral reforça que trajetórias de incapacidade são dinâmicas e dependem de fatores contextuais, terapêuticos e ocupacionais \parencite{przysada2019anassessmentofworkabilityafterthecompletionofre, deveau2011workplaceaccommodationandauditbasedevaluationprocess}. Assim, qualquer sistema decisório que privilegie exclusivamente o documento médico corre o risco de subestimar a dimensão relacional entre condição de saúde, ambiente de trabalho e possibilidade concreta de desempenho funcional.

No plano administrativo, a adoção do ATESTMED deve ser compreendida como resposta institucional a um dilema clássico de políticas públicas, que consiste em decidir em massa sem degradar a qualidade da decisão individual. A escala administrativa é imperativa em sistemas previdenciários de cobertura ampla; todavia, a expansão de capacidade por meio de simplificação procedimental pode gerar dilemas não triviais. Se, por um lado, a via documental digital assíncrona reduz deslocamentos, acelera respostas e potencialmente racionaliza o uso da força de trabalho pericial, por outro, ela pode aumentar a dependência de documentos heterogêneos, emitidos em contextos assistenciais diversos e nem sempre orientados pelos requisitos previdenciários. Estudos descritivos sobre documentação médico-administrativa de deficiência e perfis de demandantes mostram precisamente a diversidade de situações clínicas, sociais e institucionais que chegam aos sistemas de reconhecimento formal \parencite{sa2020medicalcertificateofdisabilityadescriptiveanalysisofav, hospital2016mentaldisabilityaretrospectivestudyofsocioclinica}. Em consequência, a promessa de padronização administrativa precisa ser confrontada com a variabilidade empírica dos casos.

É nesse ponto que a decisão baseada em evidências se torna eixo analítico indispensável. No âmbito deste artigo, tal expressão não é empregada em sentido retórico, como sinônimo genérico de modernização, mas em acepção metodológica, segundo a qual decisões públicas devem ser informadas por síntese rigorosa da literatura, por explicitação de critérios de inclusão e exclusão de estudos, por avaliação crítica da transferibilidade dos achados e por prudência na inferência normativa. Para sustentar essa análise, o artigo adota uma revisão estruturada da literatura, orientada por princípios PRISMA e pela recuperação, verificação e análise de textos integrais, restringindo a síntese final a estudos cuja versão integral original foi efetivamente recuperada. Essa opção busca ampliar a rastreabilidade da evidência utilizada e evitar inferências baseadas apenas em resumos ou metadados, aspecto especialmente relevante em um tema no qual nuances institucionais, clínicas e contextuais são decisivas.

A centralidade do estudo de \textcite{soares20241031t0000000000asistematicadoatestmednaconcessa} entre os estudos incluídos confere densidade empírica específica ao debate sobre o ATESTMED, sobretudo por tratar diretamente de sua sistemática e de seus impactos nos grandes números da Previdência Social. Contudo, a interpretação desse caso requer diálogo com literatura mais ampla sobre telemedicina, incapacidade, documentação médico-administrativa, avaliação funcional, incentivos institucionais e inclusão digital. Esse diálogo não pressupõe equivalência entre ATESTMED e teleperícia; ele mobiliza evidências adjacentes para compreender limites da mediação digital, desenho de fluxos híbridos e condições institucionais de uso de documentos médicos em decisões públicas. O objetivo não é transplantar mecanicamente achados internacionais para a realidade brasileira, mas construir um quadro analítico capaz de situar o ATESTMED como inovação administrativa inserida em tendências globais de digitalização, sem perder de vista as especificidades normativas e operacionais da Perícia Médica Federal.

Diante desse cenário, este artigo pergunta em que medida o ATESTMED, enquanto instrumento de saúde digital aplicado à concessão de benefício por incapacidade temporária, contribui para ampliar a eficiência e o acesso na Perícia Médica Federal sem comprometer a qualidade decisória, a consistência técnico-probatória e a aderência a uma lógica de decisão baseada em evidências. Essa pergunta desdobra-se em uma tensão analítica fundamental. De um lado, há a necessidade de respostas administrativas céleres e escaláveis em um sistema pressionado por grandes volumes de demanda. De outro, há a exigência de preservar a integridade da decisão pericial, evitando que a simplificação procedimental converta a prova médica em mera formalidade documental ou que a busca por produtividade obscureça casos que demandam exame mais aprofundado. A hipótese subjacente é que o valor público do ATESTMED não pode ser aferido apenas por indicadores de vazão administrativa, mas deve ser examinado à luz de critérios de suficiência informacional, equidade de acesso, governança da evidência e capacidade institucional de discriminar adequadamente entre casos aptos à decisão documental assíncrona e casos que requerem teleperícia ou avaliação presencial.

Assim, esta introdução sustenta que o debate sobre ATESTMED ultrapassa a dicotomia simplista entre inovação e resistência corporativa. O que está em jogo é a arquitetura de um regime decisório híbrido, no qual documentos médicos, plataformas digitais, protocolos administrativos e julgamento pericial interagem para produzir efeitos distributivos concretos sobre milhões de segurados. Em tal regime, a saúde digital pode ser vetor de democratização do acesso, mas também de opacificação de critérios; pode reduzir filas, mas também deslocar desigualdades; pode padronizar procedimentos, mas também intensificar erros de classificação se não houver base empírica suficiente. Por isso, examinar o ATESTMED a partir de uma revisão de literatura baseada na análise de textos integrais recuperados e foco em decisão baseada em evidências não é apenas uma escolha acadêmica, mas uma exigência analítica compatível com a magnitude do problema público e com a responsabilidade institucional da Perícia Médica Federal.

Este artigo não tem por objetivo analisar o ATESTMED apenas como política pública previdenciária ou inovação administrativa, mas examinar como evidências são produzidas, selecionadas, validadas e convertidas em decisão pública no âmbito da Perícia Médica Federal.

\section{Referencial analítico: decisão baseada em evidências e análise de políticas públicas}
\label{sec:orgdbc8bfe}
A análise do ATESTMED, enquanto inovação de saúde digital aplicada à concessão de benefício por incapacidade temporária no âmbito da Perícia Médica Federal, requer um referencial analítico capaz de articular, de modo simultâneo, dimensões epistemológicas, organizacionais e político-institucionais da decisão pública. Não se trata apenas de examinar se uma tecnologia digital acelera fluxos administrativos ou reduz filas; trata-se de compreender em que medida ela reconfigura os critérios de produção da verdade administrativa, os mecanismos de legitimação da decisão pericial, os padrões de eficiência e efetividade da política e, sobretudo, o desenho institucional por meio do qual evidências são selecionadas, processadas e convertidas em decisões com efeitos distributivos concretos. Nesse sentido, a decisão baseada em evidências e a análise de políticas públicas constituem eixos complementares para interpretar o ATESTMED não como ferramenta neutra, mas como arranjo decisório inserido em uma política pública de proteção social, mediada por infraestrutura digital, regras procedimentais e capacidades estatais.

\subsection{Modelo analítico-decisório adaptado de Secchi}
\label{sec:org01c6030}

Este artigo adota, como estrutura analítica complementar, uma adaptação da abordagem racionalista de análise de políticas públicas proposta por Secchi. Embora o autor não formule um modelo específico de “decisão baseada em evidências” nos termos empregados pela literatura contemporânea de DBE, sua concepção de análise de políticas públicas é diretamente compatível com esse campo, pois compreende a boa decisão pública como aquela sustentada por informações e análises confiáveis, orientada por valores socialmente aceitos e voltada à produção de efeitos desejados sobre o bem-estar coletivo. Nessa perspectiva, a análise de políticas públicas opera como método de qualificação da decisão pública, articulando diagnóstico do problema, análise de soluções, definição de critérios e indicadores, projeção de resultados e recomendação \parencite{secchi2020analisedepoliticaspublicas}.

Aplicado ao ATESTMED, esse modelo permite organizar a evidência revisada em uma sequência decisória. Primeiro, delimita-se o problema público associado à fila pericial, à demanda reprimida, à desigualdade territorial e ao risco de degradação da qualidade decisória. Em seguida, identificam-se alternativas institucionais possíveis, distinguindo análise documental assíncrona pelo ATESTMED, teleperícia síncrona, perícia presencial e modelos híbridos. Posteriormente, essas alternativas são examinadas à luz de critérios de eficiência, efetividade, equidade, suficiência probatória, auditabilidade e segurança decisória. Por fim, formula-se uma recomendação de redesenho do fluxo pericial compatível com a evidência disponível. Assim, a decisão baseada em evidências não é tratada apenas como uso de literatura científica, mas como processo estruturado de conversão de evidências em recomendação pública tecnicamente consistente, socialmente sensível e institucionalmente viável.

É nesse ponto que a análise de políticas públicas oferece instrumentos decisivos. Em vez de tratar o ATESTMED como solução técnica isolada, a abordagem de políticas públicas o insere no ciclo mais amplo de formulação, implementação, monitoramento e revisão de uma política de reconhecimento de incapacidade laboral. O foco desloca-se, então, para o desenho institucional, isto é, para quem decide, com base em quais informações, segundo quais regras, com quais incentivos, sob quais mecanismos de controle e com quais possibilidades de revisão. A literatura demonstra que, em políticas de incapacidade e deficiência, pequenas alterações procedimentais podem produzir efeitos substantivos sobre inclusão, exclusão e trajetórias de acesso a direitos \parencite{nogueira20200731t0100000100aavaliacaodadeficiencianatraje,rajasulochana2026doeshavingadisabilitycertificateensurebenefit}. Em outras palavras, a arquitetura decisória não é acessória; ela é constitutiva do próprio conteúdo da política.

\subsection{Evidência, prova documental e legitimidade da decisão pericial}
\label{sec:org6dcf73e}

Para evitar ambiguidade conceitual, este artigo distingue três modalidades de instrução e decisão pericial. O ATESTMED é compreendido como análise documental digital assíncrona, baseada em atestados, relatórios, laudos e demais documentos médicos enviados pelo segurado. A teleperícia corresponde à avaliação remota síncrona, com interação entre perito e segurado por meio tecnológico, podendo permitir entrevista, esclarecimento de dúvidas e alguma observação funcional. A perícia presencial, por sua vez, envolve avaliação direta, com possibilidade de exame físico e interação clínica plena. Assim, a literatura sobre telemedicina e avaliação remota é mobilizada de forma indireta e complementar, para discutir limites da mediação digital e justificar o escalonamento entre modalidades, e não como evidência direta de validade do ATESTMED.

No plano metodológico, este artigo adota uma compreensão ampliada de decisão baseada em evidências, afastando-se de uma leitura estritamente positivista segundo a qual a evidência seria um dado objetivo, autoevidente e automaticamente conversível em decisão correta. Em políticas públicas, especialmente naquelas que envolvem incapacidade laboral, deficiência, documentação médico-administrativa e acesso a benefícios, a evidência é sempre produzida em contextos institucionais específicos, por meio de classificações, formulários, protocolos, linguagens profissionais e critérios de admissibilidade documental. A literatura sobre incapacidade para o trabalho mostra que a avaliação não é mera leitura passiva de um estado clínico, mas processo interpretativo dependente de instrumentos, formação profissional, rotinas organizacionais e enquadramentos normativos \parencite{steenbeek2011thedevelopmentofinstrumentstomeasuretheworkdisa,dijkstrasdatamakingupincapacityforwork}. Assim, a decisão baseada em evidências, no contexto da Perícia Médica Federal, deve ser entendida como decisão informada por evidências institucionalmente curadas, hierarquizadas e operacionalizadas, e não como simples automatização de fatos médicos.

Essa formulação é particularmente relevante para o ATESTMED porque o dispositivo desloca parte do centro decisório da interação pericial presencial para a análise documental digital assíncrona, fundada em atestados, exames, relatórios e registros submetidos digitalmente. Tal deslocamento altera a ecologia da prova administrativa. Se, no modelo presencial, a observação clínica direta e a entrevista pericial compunham elementos centrais da formação do juízo, no modelo digital-documental a qualidade, completude, legibilidade e padronização dos documentos passam a ter peso ainda maior. Em consequência, a robustez da decisão depende não apenas da existência de documentos, mas da qualidade do desenho institucional que regula sua recepção, triagem, validação e interpretação. O estudo de \textcite{soares20241031t0000000000asistematicadoatestmednaconcessa} evidencia precisamente que a inovação deve ser lida à luz de seus impactos “nos grandes números” da Previdência Social, isto é, em sua capacidade de alterar escala, tempo de processamento e dinâmica de concessão. Contudo, a análise de grandes volumes decisórios não pode prescindir da indagação sobre a qualidade epistêmica e a legitimidade procedimental das decisões documentais produzidas em massa.

A literatura de decisão baseada em evidências em saúde e em políticas públicas sugere que a qualidade da decisão não decorre exclusivamente da força interna da evidência, mas também da transparência do processo de seleção, da contestabilidade dos critérios e da adequação entre tipo de evidência e problema decisório. A experiência sintetizada por \textcite{gabbay2011niceguidanceonlongtermsicknessandincapacity} é especialmente elucidativa. Ao discutir a elaboração de diretrizes sobre afastamento prolongado e incapacidade, os autores mostram que, em contextos complexos, a evidência relevante não se limita a ensaios clínicos randomizados, podendo incluir estudos observacionais, sínteses descritivas, análises de custo-efetividade e testemunho especializado, desde que submetidos a procedimentos transparentes e contestáveis. Tal perspectiva é analiticamente fecunda para o caso do ATESTMED, pois a política não pode ser avaliada apenas por indicadores de produtividade ou por métricas de tempo de resposta; ela exige um quadro de evidências plural, que inclua resultados administrativos, efeitos distributivos, barreiras de acesso, riscos de erro, padrões de equidade e impactos sobre a confiança institucional.

A legitimidade, nesse contexto, deve ser concebida em ao menos três planos articulados. O primeiro é a legitimidade legal-formal, referente à aderência do procedimento às normas que regem a concessão do benefício e a atuação da Perícia Médica Federal. O segundo é a legitimidade técnico-cognitiva, relacionada à suficiência e adequação das evidências mobilizadas para sustentar o juízo sobre incapacidade. O terceiro é a legitimidade sociopolítica, vinculada à percepção de justiça, inteligibilidade e confiabilidade do processo por parte dos usuários e da sociedade. Em políticas mediadas por plataformas digitais, esses três planos podem entrar em tensão. Um procedimento pode ser formalmente válido e administrativamente célere, mas produzir opacidade para o requerente, ampliar barreiras de acesso ou reduzir a possibilidade de apresentação contextualizada do caso. A literatura sobre telemedicina durante e após a pandemia mostra que ganhos de acesso e continuidade do cuidado coexistem com limitações importantes relativas à avaliação física, à literacia digital e à desigualdade de uso \parencite{ftouni2022challengesoftelemedicineduringthecovid19pandemica,albasheer2026telemedicineutilizationanddigitalinclusionamongpe}. Tais achados são úteis, com as devidas cautelas, para compreender o ATESTMED, pois a digitalização pode ampliar capilaridade e rapidez, mas também redistribuir desigualmente as condições de produção da prova.

Essa adaptação também exige compreender a evidência como material institucionalmente produzido, selecionado e submetido a critérios de validade. No caso deste artigo, a curadoria dos estudos analisados opera como analogia metodológica da própria decisão pericial. Assim como a pesquisa deve explicitar quais evidências foram consideradas, com que critérios e sob quais limites, a decisão administrativa baseada em documentos médicos deve preservar rastreabilidade, integridade documental e possibilidade de escrutínio. Em ambos os planos, a legitimidade do juízo depende da transparência sobre o material efetivamente mobilizado.

A dimensão institucional da capacidade pericial também merece destaque. A literatura internacional sobre médicos que atuam em avaliação securitária e de incapacidade demonstra variações relevantes de formação, especialização e instrumentos de mensuração do comportamento avaliativo \parencite{steenbeek2011thedevelopmentofinstrumentstomeasuretheworkdisa}. Isso indica que a decisão não é produzida apenas por regras abstratas ou por documentos, mas por agentes dotados de repertórios profissionais, heurísticas e margens interpretativas. Em um sistema digital como o ATESTMED, parte dessas margens pode ser comprimida por formulários, campos estruturados e protocolos automatizados; outra parte, porém, pode ser deslocada para novas zonas de discricionariedade, como a suficiência do documento anexado, a coerência entre CID e tempo de afastamento, ou a necessidade de encaminhamento para avaliação presencial. O desenho institucional adequado não elimina a discricionariedade; ele a organiza, documenta e submete a critérios de consistência e revisão.

\subsection{Critérios analíticos: eficiência, efetividade, equidade e auditabilidade}
\label{sec:org7a5b968}

A noção de eficiência, por sua vez, não pode ser reduzida ao volume administrativo processado. Em políticas públicas, eficiência envolve a relação entre recursos mobilizados, processos adotados e resultados produzidos, mas deve ser qualificada pela natureza do bem público em questão. Na concessão de benefícios por incapacidade, a busca de eficiência estritamente quantitativa pode induzir simplificações indevidas, padronizações excessivas ou incentivos à decisão orientada por metas de fluxo, e não pela suficiência probatória do caso concreto. A advertência de \textcite{kelman2009performanceimprovementandperformancedysfunctionanemp} sobre os efeitos disfuncionais de metas de desempenho é particularmente pertinente, pois indicadores podem melhorar formalmente enquanto a organização desloca comportamentos para satisfazer a métrica, produzindo distorções substantivas. Aplicado ao ATESTMED, isso significa que redução de filas, aumento de concessões processadas ou diminuição do tempo médio de análise são indicadores relevantes, porém insuficientes. Sem monitoramento de qualidade decisória, reversões, judicialização, heterogeneidade territorial e perfil dos excluídos, a eficiência pode converter-se em eficiência aparente.

A efetividade, de modo correlato, deve ser definida como a capacidade da política de produzir os resultados substantivos que justificam sua existência. No caso em exame, isso implica assegurar acesso tempestivo ao benefício para segurados efetivamente incapacitados, sem comprometer a consistência técnica da decisão nem ampliar erros de inclusão ou exclusão. A literatura sobre retorno ao trabalho, incapacidade prolongada e intervenções precoces sugere que políticas bem desenhadas precisam considerar não apenas o evento administrativo da concessão, mas a trajetória funcional do indivíduo, sua condição de trabalho e os suportes disponíveis \parencite{gabbay2011niceguidanceonlongtermsicknessandincapacity,przysada2019anassessmentofworkabilityafterthecompletionofre}. Essa observação é importante porque o ATESTMED, ao operar sobre documentação médica, tende a privilegiar a dimensão documental do problema. Entretanto, incapacidade laboral é fenômeno multidimensional, situado na interface entre condição clínica, exigências ocupacionais, contexto social e possibilidades de reabilitação. Quanto mais a política se concentra em documentos descontextualizados, maior o risco de reduzir a efetividade substantiva da decisão, ainda que aumente sua velocidade operacional.

A literatura sobre deficiência e acesso a serviços reforça essa necessidade de abordagem multidimensional. Estudos sobre barreiras de acesso à saúde e à documentação médico-administrativa mostram que a existência formal de documento médico ou administrativo não garante, por si só, fruição efetiva de direitos \parencite{m2024astudyofdisabilityseveritybarriersandfacilitatingfactor,rajasulochana2026doeshavingadisabilitycertificateensurebenefit}. Do mesmo modo, análises sobre avaliação da deficiência no benefício assistencial brasileiro indicam que os instrumentos de avaliação incorporam concepções de deficiência e modelos de proteção social que afetam diretamente o alcance redistributivo da política \parencite{nogueira20200731t0100000100aavaliacaodadeficiencianatraje}. Para o ATESTMED, isso significa que a decisão baseada em evidências não pode ser confundida com decisão baseada exclusivamente em documentos médicos. A evidência relevante, em sentido analítico, inclui também informação sobre barreiras de acesso digital, capacidade de submissão documental, desigualdades regionais e adequação dos formulários às diferentes condições de saúde.

Além disso, a literatura sobre telemedicina e avaliação digital sugere, por analogia indireta, que a mediação tecnológica não é universalmente neutra entre grupos sociais. Barreiras de conectividade, letramento digital, acessibilidade e suporte familiar influenciam a capacidade de uso das plataformas \parencite{ftouni2022challengesoftelemedicineduringthecovid19pandemica,albasheer2026telemedicineutilizationanddigitalinclusionamongpe}. Em políticas de incapacidade, tais barreiras podem ser ainda mais agudas, dado que parte dos requerentes apresenta limitações funcionais, sofrimento mental, baixa renda ou inserção precária no mercado de trabalho. A análise de políticas públicas, portanto, exige incorporar a equidade como critério de avaliação da decisão baseada em evidências. Uma decisão documental pode ser tecnicamente consistente para os casos que chegam completos ao sistema e, ainda assim, ser iníqua se o desenho institucional filtra negativamente os mais vulneráveis antes mesmo da análise de mérito.

Por fim, o referencial analítico aqui adotado sustenta que a avaliação do ATESTMED deve combinar quatro perguntas centrais. A primeira indaga quais documentos entram no processo, como são validados e quão adequados são para inferir incapacidade. A segunda examina o grau de transparência, inteligibilidade, revisabilidade e confiabilidade da decisão para usuários e instituições. A terceira investiga se a política reduz custos de transação e tempo de resposta sem degradar a qualidade decisória. A quarta avalia se o arranjo institucional amplia acesso tempestivo e justo à proteção social, preservando consistência técnica e equidade. Essas perguntas convergem no problema do desenho institucional, entendido como a configuração de regras, capacidades, tecnologias e mecanismos de controle que transforma informação em decisão pública.

Em síntese, a decisão baseada em evidências, aplicada ao ATESTMED e à Perícia Médica Federal, não pode ser concebida como simples substituição do encontro presencial por processamento digital de atestados. Trata-se de uma reconfiguração institucional da produção da decisão pericial, na qual a análise documental assíncrona deve ser diferenciada da teleperícia síncrona e da perícia presencial. A análise de políticas públicas permite demonstrar que a evidência é inseparável de seus modos de produção, curadoria e uso; que a tecnologia altera a distribuição de capacidades e barreiras; e que a qualidade da decisão depende menos da digitalização em si do que da arquitetura institucional que a sustenta. É a partir desse enquadramento que o ATESTMED pode ser examinado criticamente como instrumento de saúde digital e de proteção social, situado na interseção entre inovação administrativa, racionalidade probatória e justiça distributiva \parencite{soares20241031t0000000000asistematicadoatestmednaconcessa, gabbay2011niceguidanceonlongtermsicknessandincapacity, ftouni2022challengesoftelemedicineduringthecovid19pandemica, kelman2009performanceimprovementandperformancedysfunctionanemp}.

\section{Método: revisão estruturada PRISMA com análise de textos integrais}
\label{sec:org77338b4}
\subsection{Estratégia de busca, critérios de inclusão e extração de dados}
\label{sec:org1e45082}

Este estudo não se caracteriza como revisão bibliográfica narrativa tradicional, mas como revisão estruturada da literatura com finalidade analítico-propositiva, voltada a examinar como evidências disponíveis podem informar o desenho institucional do ATESTMED no âmbito da Perícia Médica Federal. O delineamento metodológico adotado foi orientado pelos princípios do \textit{Preferred Reporting Items for Systematic Reviews and Meta-Analyses} (PRISMA), com uma adaptação procedimental central, pela qual somente estudos cuja versão integral original foi efetivamente recuperada poderiam ingressar na síntese final. Essa opção decorre da natureza do problema de pesquisa, situado na interseção entre ATESTMED, saúde digital, telemedicina, documentação médico-pericial, incapacidade laborativa, inclusão digital e decisão baseada em evidências. Em tal campo, a leitura restrita a título, resumo ou metadados é insuficiente para apreender variáveis decisivas, como desenho institucional da intervenção, critérios de elegibilidade, arquitetura decisória, limitações operacionais, vieses de implementação, pressupostos normativos e implicações para a governança pública.

A revisão partiu de um universo preliminar de registros já estruturado para o tema da Perícia Médica Federal e de decisões baseadas em evidências, no qual foram mobilizadas bases bibliográficas de alta cobertura e mecanismos complementares de recuperação documental. Entre as fontes efetivamente habilitadas no processo, destacaram-se Elsevier Scopus e CORE, utilizadas como infraestruturas de busca e recuperação em articulação com procedimentos de verificação manual. A escolha dessas bases atendeu a dois critérios metodológicos. O primeiro foi a amplitude temática, necessária para capturar literatura dispersa entre saúde pública, telemedicina, administração pública, incapacidade para o trabalho, documentação médico-administrativa de deficiência e avaliação funcional. O segundo foi a capacidade de apoiar a recuperação de documentos integrais, condição indispensável para a regra de inclusão estabelecida. Em termos operacionais, a revisão combinou busca estruturada, inspeção de metadados, verificação de duplicidades, rastreamento de referências e busca profunda para títulos considerados prioritários.

A pergunta de revisão investigou quais evidências, modelos analíticos e achados empíricos, disponíveis em estudos com versão integral recuperada, contribuem para compreender o ATESTMED e, de modo mais amplo, os efeitos, limites e condições de validade de soluções digitais aplicadas à documentação médico-administrativa, à avaliação de incapacidade, ao acesso a benefícios e à decisão médico-administrativa. Essa pergunta exigiu um escopo suficientemente amplo para abarcar estudos diretamente centrados no ATESTMED, como o trabalho de \textcite{soares20241031t0000000000asistematicadoatestmednaconcessa}, e estudos adjacentes capazes de iluminar dimensões estruturais do fenômeno, como desafios da telemedicina e da avaliação remota \parencite{ftouni2022challengesoftelemedicineduringthecovid19pandemica}, diretrizes sobre incapacidade e afastamento prolongado \parencite{gabbay2011niceguidanceonlongtermsicknessandincapacity}, comportamento avaliativo de médicos peritos ou seguradores \parencite{steenbeek2011thedevelopmentofinstrumentstomeasuretheworkdisa}, barreiras de acesso em populações com deficiência \parencite{m2024astudyofdisabilityseveritybarriersandfacilitatingfactor, albasheer2026telemedicineutilizationanddigitalinclusionamongpe}, e estudos sobre documentos médicos, benefícios e trajetórias institucionais \parencite{rajasulochana2026doeshavingadisabilitycertificateensurebenefit, sa2020medicalcertificateofdisabilityadescriptiveanalysisofav, nogueira20200731t0100000100aavaliacaodadeficiencianatraje}. Os estudos sobre telemedicina e teleperícia foram tratados como evidência adjacente ou indireta, e não como demonstração direta da validade do ATESTMED, justamente porque análise documental assíncrona, avaliação remota síncrona e perícia presencial constituem modalidades distintas.

A noção de busca profunda designa o conjunto de procedimentos adicionais de rastreamento documental aplicados aos títulos classificados como prioritários, especialmente aqueles com alta aderência temática ao ATESTMED, à perícia ou à documentação médico-pericial da incapacidade. Foram utilizados busca por DOI, consulta cruzada em múltiplas bases, verificação em repositórios institucionais, pesquisa por versões autorais, inspeção de referências citantes e citadas, além de conferência manual de páginas de periódicos e repositórios acadêmicos. Esse procedimento foi particularmente importante para assegurar a inclusão do estudo diretamente relacionado ao ATESTMED \parencite{soares20241031t0000000000asistematicadoatestmednaconcessa}, cuja centralidade temática justificava esforço máximo de recuperação.

A etapa de elegibilidade foi conduzida sob regime de leitura integral. Os documentos recuperados foram examinados para verificar se, além da aderência temática aparente, continham densidade empírica e analítica compatível com os objetivos do artigo. Os critérios de inclusão final contemplaram disponibilidade da versão integral original, pertinência direta ou indireta robusta ao eixo ATESTMED--saúde digital--incapacidade--decisão baseada em evidências, presença de conteúdo empírico, normativo-analítico ou de revisão com utilidade para inferência institucional e possibilidade de extração de categorias relevantes para a discussão da Perícia Médica Federal. Foram admitidos estudos primários, revisões sistemáticas, revisões de escopo, análises descritivas, estudos retrospectivos, estudos metodológicos e trabalhos de natureza institucional ou qualitativa, desde que contribuíssem para a compreensão do problema decisório. Foram excluídos textos sem recuperação integral, documentos com aderência apenas nominal ao tema, estudos excessivamente periféricos e materiais cuja integridade documental não pudesse ser assegurada.

A recuperação, verificação e análise dos textos integrais constituíram o núcleo distintivo do método. Cada documento recuperado foi submetido à conferência de autoria, título, ano, DOI quando disponível, paginação detectada e compatibilidade entre metadados e conteúdo interno. Essa etapa foi necessária porque, em revisões sobre saúde digital e incapacidade, podem circular registros incompletos, versões preliminares, duplicatas imperfeitas ou documentos com metadados inconsistentes. O procedimento permitiu consolidar um conjunto final de estudos homogêneo em termos de rastreabilidade documental, ainda que heterogêneo quanto aos desenhos de estudo. O conjunto final incluiu estudos sobre telemedicina e desafios operacionais \parencite{ftouni2022challengesoftelemedicineduringthecovid19pandemica}; diretrizes para afastamento prolongado e incapacidade \parencite{gabbay2011niceguidanceonlongtermsicknessandincapacity}; inclusão digital de pessoas com deficiência no uso de telemedicina \parencite{albasheer2026telemedicineutilizationanddigitalinclusionamongpe}; barreiras de acesso à saúde entre adultos com deficiência \parencite{m2024astudyofdisabilityseveritybarriersandfacilitatingfactor}; instrumentos para mensuração do comportamento avaliativo de médicos em contextos de incapacidade laboral \parencite{steenbeek2011thedevelopmentofinstrumentstomeasuretheworkdisa}; e análises sobre documentação médico-administrativa, benefícios e trajetórias de reconhecimento institucional da deficiência \parencite{rajasulochana2026doeshavingadisabilitycertificateensurebenefit, sa2020medicalcertificateofdisabilityadescriptiveanalysisofav, nogueira20200731t0100000100aavaliacaodadeficiencianatraje}.

A assistência por inteligência artificial foi empregada em caráter instrumental e subordinado à supervisão metodológica humana. Seu uso incidiu sobre tarefas de apoio à curadoria, como normalização de metadados, sugestão de agrupamentos temáticos preliminares, identificação de possíveis duplicidades, extração de elementos bibliográficos recorrentes, localização de trechos potencialmente relevantes nos textos integrais e construção de matriz de leitura comparativa. A IA não operou como instância autônoma de julgamento científico. Toda sugestão automatizada foi submetida a verificação humana, e toda decisão de elegibilidade, exclusão ou inclusão foi tomada por leitura crítica do texto integral.

A extração de dados foi realizada a partir da leitura integral dos 20 estudos incluídos, com registro em matriz analítica própria de informações relativas a autoria, ano, país ou contexto institucional, tipo de estudo, objeto principal, desenho metodológico, população ou unidade de análise, conceitos centrais, achados empíricos, limitações declaradas e contribuição potencial para o problema do ATESTMED. Também foram extraídas categorias transversais relevantes para a síntese interpretativa, incluindo mediação tecnológica da decisão, qualidade e suficiência documental, barreiras de acesso, assimetrias de inclusão digital, critérios de avaliação funcional, efeitos de desenho institucional, riscos de simplificação decisória e relações entre evidência, padronização e discricionariedade profissional. Essa estratégia permitiu comparar estudos de natureza distinta sem forçar homogeneidade indevida, preservando a pluralidade metodológica dos estudos incluídos.

A síntese adotada foi predominantemente qualitativa, de caráter narrativo-analítico, e não metanalítico. Essa escolha decorreu da heterogeneidade dos desenhos, dos desfechos e das unidades de observação presentes nos estudos incluídos. Havia, por exemplo, revisões sistemáticas sobre telemedicina \parencite{ftouni2022challengesoftelemedicineduringthecovid19pandemica}, diretrizes baseadas em múltiplos tipos de evidência \parencite{gabbay2011niceguidanceonlongtermsicknessandincapacity}, estudos metodológicos sobre comportamento avaliativo de médicos \parencite{steenbeek2011thedevelopmentofinstrumentstomeasuretheworkdisa}, análises descritivas de documentação médico-administrativa de deficiência \parencite{sa2020medicalcertificateofdisabilityadescriptiveanalysisofav}, estudos retrospectivos sobre perfis clínico-sociais \parencite{hospital2016mentaldisabilityaretrospectivestudyofsocioclinica}, investigações sobre acesso geográfico a reabilitação \parencite{gao2019disabilityconcentrationandaccesstorehabilitationservice} e trabalhos sobre incentivos, benefícios e trajetórias institucionais \parencite{corden2001incapacitybenefitsandworkincentives, sainsbury2006routesontoincapacitybenefitsfindingsfromqualitati}. Em tal contexto, a comparabilidade estatística entre efeitos substantivos seria artificial e metodologicamente frágil. A estatística foi utilizada apenas em chave descritiva, para tornar transparente o fluxo de seleção e a regra de parada da revisão.

\subsection{Descrição quantitativa do fluxo de seleção}
\label{sec:org9c490cf}

Além da síntese qualitativa, realizou-se análise estatística descritiva simples do fluxo de seleção dos estudos, com cálculo de frequências absolutas e relativas. Do total de 248 registros lidos no conjunto preliminar, 20 estudos foram incluídos na síntese final, correspondendo a 8,1\% dos registros inicialmente avaliados. Outros 12 registros foram excluídos pela impossibilidade de recuperação da versão integral, o equivalente a 4,8\% do conjunto preliminar. Considerando apenas os 32 registros que chegaram à etapa de recuperação ou verificação da versão integral, a taxa de inclusão foi de 62,5\%, enquanto a taxa de exclusão por indisponibilidade documental foi de 37,5\%. Por fim, 216 registros, correspondentes a 87,1\% do conjunto preliminar, deixaram de ser avaliados após o atingimento do alvo amostral previamente definido. Ao final da busca profunda, não restaram títulos prioritários sem versão integral recuperada.

\begin{table}[htbp]
\caption{Descrição quantitativa do fluxo de seleção dos estudos}
\centering
\footnotesize
\begin{tabular}{@{}p{0.66\textwidth}rr@{}}
Etapa & n & \%\\
\hline
Registros lidos no conjunto preliminar & 248 & 100,0\\
Estudos incluídos na síntese final & 20 & 8,1\\
Registros excluídos por indisponibilidade da versão integral & 12 & 4,8\\
Registros não avaliados após atingimento do alvo amostral & 216 & 87,1\\
\end{tabular}
\end{table}

Fonte: elaboração própria.

\begin{table}[htbp]
\caption{Taxa de recuperação entre registros submetidos à verificação da versão integral}
\centering
\footnotesize
\begin{tabular}{@{}p{0.66\textwidth}rr@{}}
Situação após verificação da versão integral & n & \%\\
\hline
Versão integral recuperada e estudo incluído & 20 & 62,5\\
Versão integral não recuperada e registro excluído & 12 & 37,5\\
Total submetido à recuperação ou verificação integral & 32 & 100,0\\
\end{tabular}
\end{table}

Fonte: elaboração própria.

Esses resultados não devem ser interpretados como estimativa de prevalência temática da literatura sobre ATESTMED, saúde digital pericial ou documentação médico-administrativa. Eles descrevem o processo de seleção adotado e cumprem função de transparência metodológica. A descrição quantitativa aumenta a auditabilidade da passagem do universo preliminar de registros para o conjunto final de estudos analisados, explicita a regra de parada e evita que o leitor confunda a síntese analítico-propositiva com uma metanálise ou com um mapeamento exaustivo de toda a produção disponível.

Por fim, importa registrar que o método foi concebido para maximizar validade interna da síntese, rastreabilidade documental e transparência decisória. A exigência de versão integral original reduziu o risco de inferências baseadas em informação incompleta; a curadoria assistida por IA ampliou a capacidade de organização e recuperação; a decisão humana final preservou o juízo crítico necessário à interpretação de um objeto institucionalmente sensível; e a estrutura PRISMA forneceu disciplina procedimental ao fluxo de identificação, triagem, elegibilidade e inclusão. Em um campo como o da Perícia Médica Federal, no qual soluções digitais como o ATESTMED podem alterar escalas de acesso, tempos de decisão, padrões de prova e formas de reconhecimento da incapacidade, a robustez metodológica da revisão não é um requisito meramente formal. Ela constitui condição substantiva para que a discussão sobre inovação administrativa permaneça ancorada em evidências efetivamente examinadas, e não em promessas tecnológicas, impressões anedóticas ou generalizações sem lastro empírico \parencite{soares20241031t0000000000asistematicadoatestmednaconcessa, ftouni2022challengesoftelemedicineduringthecovid19pandemica, gabbay2011niceguidanceonlongtermsicknessandincapacity}.

\section{Resultados da revisão estruturada}
\label{sec:org4ae4ef0}
Os resultados são apresentados em três componentes, compreendendo uma matriz textual dos 20 estudos incluídos, a síntese dos achados por eixo temático e uma síntese integradora voltada à formulação da recomendação analítico-propositiva.

\subsection{Matriz dos estudos incluídos}
\label{sec:orgec3140d}

A matriz abaixo consolida, em formato textual, os 20 estudos com versão integral recuperada, evitando tabela extensa de difícil leitura e preservando a rastreabilidade das citações.

\begin{itemize}
\item 1. A sistemática do ATESTMED na concessão do atestado por incapacidade e seu impacto nos grandes números da Previdência Social (2024). Eixo nacional sobre ATESTMED, concessão documental e efeitos institucionais. \parencite{soares20241031t0000000000asistematicadoatestmednaconcessa}.
\item 2. \textit{Challenges of Telemedicine during the COVID-19 pandemic: a systematic review} (2022). Evidência sobre telemedicina, consulta remota ou mediação digital da avaliação. \parencite{ftouni2022challengesoftelemedicineduringthecovid19pandemica}.
\item 3. \textit{NICE guidance on long-term sickness and incapacity} (2011). Evidência sobre avaliação de incapacidade, funcionalidade e documentação médica. \parencite{gabbay2011niceguidanceonlongtermsicknessandincapacity}.
\item 4. \textit{Telemedicine utilization and digital inclusion among people with disabilities in Saudi Arabia} (2026). Evidência sobre telemedicina, consulta remota ou mediação digital da avaliação. \parencite{albasheer2026telemedicineutilizationanddigitalinclusionamongpe}.
\item 5. \textit{A Study of Disability Severity, Barriers, and Facilitating Factors in Accessing Healthcare Among Differently Abled Adults} (2024). Evidência sobre barreiras de acesso, deficiência e serviços de saúde. \parencite{m2024astudyofdisabilityseveritybarriersandfacilitatingfactor}.
\item 6. \textit{The development of instruments to measure the work disability assessment behaviour of insurance physicians} (2011). Evidência sobre comportamento avaliativo de médicos peritos e padronização decisória. \parencite{steenbeek2011thedevelopmentofinstrumentstomeasuretheworkdisa}.
\item 7. \textit{Does Having a Disability Certificate Ensure Benefits?} (2026). Evidência sobre reconhecimento documental de deficiência e acesso efetivo a benefícios. \parencite{rajasulochana2026doeshavingadisabilitycertificateensurebenefit}.
\item 8. \textit{Making up incapacity for work?} (s.d.). Evidência teórica sobre construção institucional da incapacidade laboral. \parencite{dijkstrasdatamakingupincapacityforwork}.
\item 9. \textit{Pediatric complex chronic conditions classification system version 2} (2014). Evidência sobre classificação, padronização e complexidade clínica. \parencite{feudtner2014pediatriccomplexchronicconditionsclassificationsyst}.
\item 10. \textit{Routes Onto Incapacity Benefits: Findings from Qualitative Research} (2006). Evidência qualitativa sobre trajetórias de entrada em benefícios por incapacidade. \parencite{sainsbury2006routesontoincapacitybenefitsfindingsfromqualitati}.
\item 11. \textit{An assessment of work ability after the completion of rehabilitation} (2019). Evidência sobre reabilitação, capacidade laboral e retorno ao trabalho. \parencite{przysada2019anassessmentofworkabilityafterthecompletionofre}.
\item 12. \textit{Medical certificate of disability: a descriptive analysis of Aveiro region, from 2010 to 2016} (2020). Evidência sobre documentos médicos de deficiência. \parencite{sa2020medicalcertificateofdisabilityadescriptiveanalysisofav}.
\item 13. \textit{Assessment of Aphasia Across the International Classification of Functioning, Disability and Health Using an iPad-Based Application} (2017). Evidência sobre avaliação funcional digital e Classificação Internacional de Funcionalidade. \parencite{guo2017assessmentofaphasiaacrosstheinternationalclassification}.
\item 14. \textit{Mental Disability: A Retrospective Study of Socio-Clinical Profile of Patients Seeking Disability Certificate at a Tertiary Care Centre in Delhi} (2016). Evidência sobre documentação médica de deficiência mental e perfil sociodemográfico. \parencite{hospital2016mentaldisabilityaretrospectivestudyofsocioclinica}.
\item 15. \textit{Disability concentration and access to rehabilitation services: a pilot spatial assessment applying geographic information system analysis} (2019). Evidência sobre barreiras territoriais e acesso à reabilitação. \parencite{gao2019disabilityconcentrationandaccesstorehabilitationservice}.
\item 16. \textit{A meta-review of literature reviews of disability, travel and inequalities} (2025). Evidência sobre deficiência, mobilidade, barreiras de acesso e desigualdades. \parencite{mackett2025ametareviewofliteraturereviewsofdisabilitytravel}.
\item 17. A avaliação da deficiência na trajetória do benefício de prestação continuada (2020). Evidência brasileira sobre avaliação da deficiência e proteção social. \parencite{nogueira20200731t0100000100aavaliacaodadeficiencianatraje}.
\item 18. \textit{Incapacity Benefits and Work Incentives} (2001). Evidência sobre benefícios por incapacidade e incentivos ao trabalho. \parencite{corden2001incapacitybenefitsandworkincentives}.
\item 19. \textit{Performance Improvement and Performance Dysfunction} (2009). Evidência sobre indicadores, metas e disfunções de desempenho. \parencite{kelman2009performanceimprovementandperformancedysfunctionanemp}.
\item 20. \textit{Workplace accommodation and audit-based evaluation process for compliance with the Employment Equity Act} (2011). Evidência sobre acomodação laboral, auditoria e inclusão. \parencite{deveau2011workplaceaccommodationandauditbasedevaluationprocess}.
\end{itemize}

\subsection{Achados por eixo temático}
\label{sec:org9aee4d1}

O primeiro resultado relevante da revisão consiste na constatação de que o ATESTMED emerge, nos estudos incluídos, menos como uma tecnologia isolada e mais como um rearranjo institucional da produção de prova médico-administrativa em larga escala. O estudo diretamente dedicado ao tema demonstra que a sistemática do ATESTMED deve ser compreendida como mecanismo de racionalização do fluxo concessório, com impacto sobre os “grandes números” da Previdência Social, especialmente ao deslocar parte da centralidade da interação pericial presencial para a consistência documental e para a padronização procedimental \parencite{soares20241031t0000000000asistematicadoatestmednaconcessa}. Esse achado é particularmente importante porque desloca o debate da dicotomia simplista entre “presencial” e “digital” para uma questão mais substantiva, referente aos atributos documentais, clínicos e funcionais necessários para que uma decisão automatizada ou semipresencial preserve validade, legitimidade e capacidade de controle. A literatura revisada sugere que o ATESTMED não pode ser avaliado apenas por sua eficiência operacional; ele deve ser examinado como dispositivo de triagem probatória, cuja robustez depende da qualidade dos atestados, relatórios e laudos médicos, da inteligibilidade do nexo entre diagnóstico e incapacidade, da temporalidade do agravo e da possibilidade de inferência funcional a partir do documento apresentado.

Nesse sentido, o eixo da avaliação documental revelou um padrão recorrente, segundo o qual documentos médicos são social e administrativamente potentes, mas epistemicamente desiguais. Estudos sobre documentação médico-administrativa de deficiência e concessão de benefícios mostram que a existência de um documento médico ou administrativo não assegura, por si só, acesso efetivo a direitos, seja por barreiras burocráticas, seja por assimetrias na interpretação institucional do conteúdo documentado \parencite{rajasulochana2026doeshavingadisabilitycertificateensurebenefit}. A análise da trajetória da avaliação da deficiência no benefício assistencial brasileiro reforça essa conclusão ao demonstrar que o documento médico opera dentro de uma cadeia decisória mais ampla, na qual categorias administrativas, critérios de elegibilidade e leituras periciais reconfiguram o valor da prova clínica \parencite{nogueira20200731t0100000100aavaliacaodadeficiencianatraje}. Em paralelo, estudos descritivos sobre documentos médicos de deficiência em contextos regionais evidenciam forte variabilidade de perfis clínicos, faixas etárias e motivos de emissão documental, indicando que a padronização formal do documento não elimina a heterogeneidade substantiva dos casos \parencite{sa2020medicalcertificateofdisabilityadescriptiveanalysisofav, hospital2016mentaldisabilityaretrospectivestudyofsocioclinica}. A revisão indica, portanto, que a avaliação documental pode sustentar decisões céleres em casos de baixa ambiguidade, mas sua confiabilidade diminui quando o documento é lacônico, quando o diagnóstico não traduz limitação funcional ou quando o contexto ocupacional é decisivo para a caracterização da incapacidade.

A literatura também mostrou que a avaliação documental é particularmente sensível ao problema da tradução entre doença, deficiência, funcionalidade e incapacidade para o trabalho. Esse ponto aparece de forma explícita em trabalhos que problematizam a própria construção social e institucional da incapacidade laboral \parencite{dijkstrasdatamakingupincapacityforwork}. A incapacidade não se apresenta como dado bruto imediatamente extraível do diagnóstico; ela é produzida por uma operação interpretativa que articula achados clínicos, exigências laborais, prognóstico, contexto social e critérios normativos. Em consequência, sistemas como o ATESTMED tendem a funcionar melhor quando o documento médico já incorpora, de forma estruturada, elementos que tradicionalmente seriam explorados em entrevista pericial, como duração estimada do afastamento, restrições funcionais objetivas, terapêutica instituída, resposta ao tratamento e justificativa clínica para a limitação laboral. A ausência desses elementos amplia a incerteza decisória e pode induzir tanto indeferimentos por insuficiência probatória quanto concessões baseadas em documentação formalmente válida, porém materialmente frágil.

No eixo da telemedicina e da avaliação remota, os resultados da revisão indicam que a digitalização do cuidado e da avaliação em saúde oferece oportunidades reais de ampliação de acesso, mas produz novas camadas de exclusão e novos riscos de qualidade. A revisão sistemática sobre os desafios da telemedicina durante a pandemia de COVID-19 identificou barreiras recorrentes relacionadas à infraestrutura tecnológica, segurança da informação, treinamento profissional, aceitabilidade pelos usuários e adequação regulatória \parencite{ftouni2022challengesoftelemedicineduringthecovid19pandemica}. Esses achados não são diretamente transferíveis ao ATESTMED, porque telemedicina e teleperícia pressupõem alguma interação remota, enquanto o ATESTMED opera predominantemente por análise documental assíncrona. Ainda assim, eles são relevantes como evidência indireta sobre os limites da mediação digital em decisões que dependem de informação clínica qualificada. O estudo sobre utilização de telemedicina entre pessoas com deficiência na Arábia Saudita aprofunda essa conclusão ao mostrar baixa adoção global do recurso, apesar de avaliações favoráveis quanto à privacidade e à facilidade de uso entre usuários. Apenas 25,4\% dos participantes relataram uso de telemedicina, e os principais obstáculos foram preferência por atendimento presencial, baixa familiaridade com a tecnologia e disponibilidade limitada de provedores \parencite{albasheer2026telemedicineutilizationanddigitalinclusionamongpe}. O dado é relevante para a Perícia Médica Federal porque sugere que a expansão de soluções digitais não pode pressupor homogeneidade de acesso, letramento digital ou capacidade de navegação burocrática por parte dos segurados.

A revisão mostrou, ademais, que a inclusão digital é variável crítica para qualquer arquitetura decisória baseada em saúde digital. Pessoas com maior escolaridade e residentes em áreas urbanas tendem a utilizar mais recursos remotos \parencite{albasheer2026telemedicineutilizationanddigitalinclusionamongpe}, ao passo que indivíduos com maiores limitações funcionais, menor escolaridade e inserção social mais vulnerável enfrentam barreiras cumulativas \parencite{m2024astudyofdisabilityseveritybarriersandfacilitatingfactor, gao2019disabilityconcentrationandaccesstorehabilitationservice, mackett2025ametareviewofliteraturereviewsofdisabilitytravel}. Em termos analíticos, isso significa que a digitalização do acesso e da instrução pericial pode aumentar eficiência média e, simultaneamente, aprofundar desigualdades marginais. Os estudos revisados convergem para a necessidade de desenho inclusivo, com canais híbridos, suporte assistido, acessibilidade comunicacional, interoperabilidade de sistemas e mecanismos de revisão humana para casos complexos. Em outras palavras, a saúde digital só se converte em instrumento de decisão baseada em evidências quando a infraestrutura sociotécnica permite que a evidência seja produzida, transmitida e interpretada sem exclusão sistemática de determinados grupos.

No eixo da avaliação funcional, a revisão produziu um dos resultados mais consistentes e teoricamente fecundos. Diversos estudos convergem para a insuficiência do diagnóstico isolado como fundamento decisório em matéria de incapacidade e deficiência. O uso de instrumentos estruturados, escalas e classificações orientadas à funcionalidade aparece como estratégia central para reduzir variabilidade, aumentar comparabilidade e aproximar a decisão administrativa da experiência concreta de limitação \parencite{steenbeek2011thedevelopmentofinstrumentstomeasuretheworkdisa, guo2017assessmentofaphasiaacrosstheinternationalclassification, feudtner2014pediatriccomplexchronicconditionsclassificationsyst}. O estudo de Steenbeek et al. é particularmente relevante porque trata do desenvolvimento de instrumentos para mensurar o comportamento avaliativo de médicos peritos de seguros, evidenciando que a própria conduta do avaliador pode e deve ser objeto de padronização e mensuração \parencite{steenbeek2011thedevelopmentofinstrumentstomeasuretheworkdisa}. Esse achado desloca o foco da variabilidade apenas do paciente para a variabilidade do processo decisório, o que é crucial em contextos de grande escala, como o da Previdência Social brasileira.

A incorporação de referenciais funcionais, próximos à Classificação Internacional de Funcionalidade, Incapacidade e Saúde, também se mostrou promissora para ambientes digitais. O estudo sobre avaliação de afasia com aplicação baseada em \textit{iPad} demonstrou a viabilidade de instrumentos digitais para captar dimensões funcionais complexas de comunicação e participação \parencite{guo2017assessmentofaphasiaacrosstheinternationalclassification}. Embora o objeto clínico seja específico, o resultado metodológico é mais amplo, indicando que plataformas digitais podem ser desenhadas não apenas para receber documentos, mas para estruturar a coleta de informação funcional de modo padronizado, reprodutível e auditável. Para o ATESTMED, isso sugere a possibilidade de evoluir de um modelo predominantemente documental para um modelo documental-funcional, no qual o atestado médico seja complementado por campos estruturados sobre limitações, atividades comprometidas, exigências ocupacionais e prognóstico. Tal desenho reduziria a distância entre prova clínica e decisão pericial, especialmente nos casos em que o diagnóstico é pouco informativo sobre a capacidade laboral.

A revisão também evidenciou que a avaliação funcional tem implicações diretas para o retorno ao trabalho. Diretrizes sobre afastamento prolongado e incapacidade enfatizam intervenções precoces, coordenação entre atores e foco na manutenção ou recuperação da capacidade laboral residual \parencite{gabbay2011niceguidanceonlongtermsicknessandincapacity}. Estudos sobre trajetórias de entrada em benefícios por incapacidade mostram que o afastamento prolongado frequentemente resulta de processos cumulativos, nos quais fatores clínicos, laborais, familiares e institucionais se reforçam mutuamente \parencite{sainsbury2006routesontoincapacitybenefitsfindingsfromqualitati, corden2001incapacitybenefitsandworkincentives}. Já a avaliação da capacidade para o trabalho após reabilitação em pessoas com doenças musculoesqueléticas crônicas indica que a mensuração funcional é decisiva para distinguir melhora clínica de efetiva prontidão ocupacional \parencite{przysada2019anassessmentofworkabilityafterthecompletionofre}. O resultado sintético desse eixo é que decisões periciais baseadas apenas em presença de doença ou em temporalidade do tratamento tendem a ser menos informativas do que aquelas ancoradas em funcionalidade, possibilidade de adaptação e potencial de reinserção.

Quanto ao eixo da inteligência artificial, a síntese final não contém estudos centrados exclusivamente em IA aplicada à perícia médica federal, mas oferece elementos suficientes para uma inferência analítica prudente. Em primeiro lugar, a literatura sobre instrumentos de avaliação, classificação e padronização sugere que a IA só pode operar de forma confiável quando alimentada por dados estruturados, semanticamente consistentes e clinicamente relevantes \parencite{steenbeek2011thedevelopmentofinstrumentstomeasuretheworkdisa, feudtner2014pediatriccomplexchronicconditionsclassificationsyst}. Em segundo lugar, os estudos sobre disfunções induzidas por metas e indicadores em sistemas públicos alertam para o risco de que tecnologias de apoio à decisão, quando subordinadas a objetivos estritamente quantitativos, produzam distorções comportamentais e degradem a qualidade substantiva da decisão \parencite{kelman2009performanceimprovementandperformancedysfunctionanemp}. Aplicado ao ATESTMED, esse resultado é particularmente sensível, uma vez que algoritmos de triagem, priorização ou deferimento automático podem aumentar produtividade, mas também incentivar preenchimento estratégico de atestados, seleção adversa de casos ou acomodação institucional a métricas de velocidade em detrimento da acurácia.

A revisão permite, assim, formular um resultado interpretativo importante, pois a IA, no contexto da Perícia Médica Federal, aparece menos como substituta da decisão pericial e mais como tecnologia de apoio à estratificação de risco documental, detecção de inconsistências, classificação de casos e priorização de revisão humana. Essa conclusão decorre da própria natureza dos achados sobre incapacidade, que é multidimensional, contextual e normativamente mediada \parencite{dijkstrasdatamakingupincapacityforwork}. Em casos simples, com documentação robusta e padrão clínico previsível, modelos algorítmicos podem contribuir para celeridade e uniformidade. Em casos complexos, contudo, a opacidade do contexto funcional, a necessidade de juízo clínico e a relevância do trabalho concreto limitam a automação decisória plena. A evidência revisada, portanto, sustenta uma arquitetura de IA assistiva, supervisionada e auditável, e não uma automação irrestrita da concessão ou do indeferimento.

No eixo da documentação médico-administrativa, os estudos incluídos revelam que atestados, relatórios, laudos e documentos equivalentes funcionam simultaneamente como instrumentos clínicos, documentos jurídico-administrativos e artefatos de política pública. Pesquisas sobre documentos médicos de deficiência em diferentes contextos mostram que sua emissão responde a padrões epidemiológicos, incentivos institucionais e expectativas sociais de acesso a benefícios \parencite{sa2020medicalcertificateofdisabilityadescriptiveanalysisofav, hospital2016mentaldisabilityaretrospectivestudyofsocioclinica, rajasulochana2026doeshavingadisabilitycertificateensurebenefit}. Em populações em idade produtiva, a busca por esse reconhecimento documental frequentemente se articula a necessidades de renda, mobilidade, inserção ocupacional e reconhecimento estatal. Isso significa que o documento médico não é mero reflexo neutro de um estado clínico; ele também pode funcionar como porta de entrada para regimes de proteção social. Para o ATESTMED, esse resultado é decisivo, porque o sistema depende da qualidade, da completude e da honestidade epistêmica dos atestados, relatórios e laudos apresentados. Se a documentação é excessivamente genérica, se não discrimina permanência, gravidade, funcionalidade e temporalidade, ela fragiliza a decisão. Se, ao contrário, é excessivamente orientada por incentivos externos, pode induzir formulações clínicas adaptadas ao benefício esperado.

A revisão sugere, por conseguinte, que os documentos médicos aptos a sustentar decisão baseada em evidências devem obedecer a requisitos de completude, padronização semântica e pertinência funcional. O estudo sobre a sistemática do ATESTMED já aponta, de forma implícita, para a centralidade dessa qualificação documental \parencite{soares20241031t0000000000asistematicadoatestmednaconcessa}. Em diálogo com a literatura internacional, torna-se evidente que a melhoria dos atestados, relatórios e laudos não depende apenas de formulário eletrônico, mas de governança clínica, com capacitação de médicos assistentes, campos obrigatórios orientados à incapacidade, validação de consistência e integração com classificações diagnósticas e funcionais. Sem isso, a digitalização apenas acelera a circulação de documentos heterogêneos.

Por fim, no eixo do retorno ao trabalho, os resultados da revisão indicam que a decisão pericial de qualidade não se esgota na definição binária entre concessão e não concessão. A literatura sobre afastamento prolongado, incentivos ao trabalho, reabilitação e acomodação laboral mostra que a incapacidade deve ser tratada como condição dinâmica, suscetível a reavaliação, adaptação e reinserção progressiva \parencite{gabbay2011niceguidanceonlongtermsicknessandincapacity, corden2001incapacitybenefitsandworkincentives, przysada2019anassessmentofworkabilityafterthecompletionofre, deveau2011workplaceaccommodationandauditbasedevaluationprocess}. Em especial, estudos sobre acomodação no local de trabalho demonstram que práticas formalmente inclusivas podem, paradoxalmente, excluir quando se limitam à conformidade procedimental e não enfrentam as barreiras reais ao desempenho \parencite{deveau2011workplaceaccommodationandauditbasedevaluationprocess}. Esse achado é relevante para a Perícia Médica Federal porque sugere que a decisão baseada em evidências deve dialogar com políticas de reabilitação, readaptação e suporte ocupacional, sob pena de transformar a perícia em mecanismo de mera classificação estática.

\subsection{Síntese integradora dos achados}
\label{sec:org05172be}

Em síntese, as evidências convergem para a ideia de que o ATESTMED representa inovação relevante na interface entre saúde digital, proteção social e gestão pericial, mas sua legitimidade técnico-científica depende de quatro condições analíticas. A primeira é a qualificação da evidência documental, com ênfase em informação funcional, temporalidade clínica e pertinência ocupacional, e não apenas em diagnóstico. A segunda é a mitigação das desigualdades de acesso digital e de produção documental, sob pena de a inovação ampliar assimetrias preexistentes. A terceira é a padronização auditável do comportamento decisório, com monitoramento da variabilidade entre avaliadores e revisão contínua dos critérios embutidos no sistema. A quarta é a articulação do processo decisório com desfechos substantivos de reabilitação, retorno ao trabalho e proteção social adequada.

Há, contudo, divergências e lacunas importantes nos estudos analisados. Parte dos estudos enfatiza a necessidade de intervenção precoce e retorno ao trabalho; outra parte alerta para o caráter socialmente construído e moralmente carregado das categorias de incapacidade \parencite{gabbay2011niceguidanceonlongtermsicknessandincapacity, dijkstrasdatamakingupincapacityforwork}. Alguns trabalhos privilegiam a eficiência classificatória e a padronização; outros mostram que classificações podem excluir experiências corporais e contextuais difíceis de codificar \parencite{steenbeek2011thedevelopmentofinstrumentstomeasuretheworkdisa, deveau2011workplaceaccommodationandauditbasedevaluationprocess}. Em vez de invalidarem a síntese, tais tensões indicam o núcleo do problema pericial contemporâneo, que consiste em decidir com celeridade sem reduzir complexidade, padronizar sem desumanizar, digitalizar sem excluir e controlar sem moralizar. Para a Perícia Médica Federal, a principal implicação é que a incorporação do ATESTMED deve ser acompanhada por desenho institucional reflexivo, avaliação contínua e produção sistemática de evidências próprias, capazes de retroalimentar critérios, fluxos e salvaguardas. Somente assim a saúde digital poderá operar, no campo pericial, não como mera automação de procedimentos, mas como infraestrutura de decisão pública efetivamente baseada em evidências.

\section{Discussão}
\label{sec:org3b7cfa9}
Os achados desta revisão permitem discutir o ATESTMED não apenas como inovação procedimental na Perícia Médica Federal, mas como dispositivo sociotécnico de reorganização da decisão administrativa em saúde digital. A discussão parte da síntese qualitativa dos estudos incluídos e se concentra nas implicações institucionais do modelo, especialmente validade da decisão documental, risco de erro, assimetria informacional, equidade digital, padronização, auditoria e capacidade estatal.

Nesse quadro, o ATESTMED emerge como resposta estatal a um problema clássico de capacidade administrativa, ligado ao processamento de grandes volumes de requerimentos por incapacidade temporária com maior celeridade, menor custo transacional e algum grau de padronização decisória. O estudo de \textcite{soares20241031t0000000000asistematicadoatestmednaconcessa} é central para essa compreensão ao situar o sistema no contexto dos “grandes números” da Previdência Social. A contribuição mais importante desse trabalho, para fins desta discussão, não reside apenas em descrever a sistemática, mas em evidenciar que a digitalização do fluxo decisório altera a natureza da prova considerada suficiente para a concessão administrativa. Em vez de uma perícia presencial como regra universal, passa-se a admitir, em hipóteses delimitadas, uma decisão fundada em documentação médica submetida digitalmente. Essa inflexão desloca o centro da decisão do exame clínico direto para a avaliação documental mediada por plataforma, com consequências imediatas para validade, risco de erro e governança.

A primeira implicação diz respeito à validade da decisão pericial. Em termos analíticos, a validade do ATESTMED não pode ser aferida apenas por indicadores de produtividade ou redução de filas; ela depende da correspondência entre a informação disponível no processo documental digital e o construto jurídico-administrativo que se pretende decidir, isto é, a incapacidade laborativa temporária nos termos normativos aplicáveis. É importante distinguir, nesse ponto, o ATESTMED da teleperícia. O primeiro corresponde predominantemente a uma análise documental assíncrona, baseada em atestados, relatórios, laudos e demais documentos médicos previamente enviados. A segunda pressupõe interação remota síncrona entre perito e segurado, podendo permitir entrevista, esclarecimento de dúvidas e alguma observação funcional, embora também apresente limitações em relação ao exame presencial.

Por essa razão, a literatura internacional sobre telemedicina e avaliação remota não deve ser tomada como evidência direta de validade do ATESTMED, mas como evidência indireta sobre os limites e as possibilidades de decisões mediadas por tecnologia. Essa literatura sugere que modalidades remotas podem ampliar acesso e continuidade, mas apresentam restrições quando a decisão depende de exame físico, observação contextual, testes funcionais ou interação clínica mais densa \parencite{ftouni2022challengesoftelemedicineduringthecovid19pandemica}. No caso específico do ATESTMED, essa cautela é ainda mais relevante, pois a decisão documental dispõe de menos elementos interacionais do que a teleperícia. Assim, o ATESTMED tende a ser mais adequado em situações clínicas com documentação robusta, padronizável, temporalmente coerente e com baixa dependência de achados semiológicos presenciais; e menos adequado em quadros nos quais a funcionalidade, a intensidade sintomática, a flutuação clínica ou a relação entre doença e trabalho exigem apreciação pericial mais densa.

Essa observação aproxima o debate brasileiro de discussões mais amplas sobre avaliação de incapacidade para o trabalho. \textcite{gabbay2011niceguidanceonlongtermsicknessandincapacity} já indicavam, no âmbito das diretrizes NICE, que decisões sobre afastamento prolongado e incapacidade não se reduzem ao diagnóstico biomédico, devendo incorporar funcionalidade, contexto ocupacional e possibilidades de retorno ao trabalho. De modo semelhante, \textcite{steenbeek2011thedevelopmentofinstrumentstomeasuretheworkdisa} demonstram que o comportamento avaliativo de médicos peritos ou seguradores não é puramente técnico, mas estruturado por instrumentos, rotinas, heurísticas e padrões institucionais. O ponto crítico, portanto, é que a digitalização não elimina a discricionariedade técnica; ela a reconfigura. No ATESTMED, a discricionariedade deixa de se apoiar predominantemente na interação presencial e passa a depender da qualidade do documento, da inteligibilidade do atestado, da consistência interna dos dados e dos filtros procedimentais da plataforma. Isso pode reduzir variabilidade arbitrária em alguns segmentos, mas também pode introduzir novas fontes de erro sistemático.

O risco de erro decisório, nesse contexto, deve ser examinado em dupla direção, abrangendo falsos positivos e falsos negativos. O debate público frequentemente enfatiza o risco de concessões indevidas, sobretudo em sistemas de análise documental. Contudo, do ponto de vista da decisão baseada em evidências e da justiça administrativa, o risco simétrico de indeferimentos indevidos é igualmente relevante. A assimetria informacional entre requerente e Estado não opera apenas como risco moral do segurado; ela também se manifesta como insuficiência epistêmica do decisor diante de documentos heterogêneos, incompletos ou produzidos em contextos assistenciais desiguais. \textcite{dijkstrasdatamakingupincapacityforwork} problematiza precisamente a construção social e institucional da incapacidade para o trabalho, sugerindo que a decisão não descobre simplesmente um fato pré-existente, mas o produz a partir de categorias, evidências e convenções. Aplicado ao ATESTMED, isso significa que a plataforma não apenas recebe informação, mas também molda o que conta como informação relevante, o que é legível, o que é comparável e o que se torna decidível.

A assimetria informacional é agravada por desigualdades de acesso e de letramento digital. A literatura sobre inclusão digital em saúde mostra que a adoção de serviços remotos é socialmente estratificada. O estudo de \textcite{albasheer2026telemedicineutilizationanddigitalinclusionamongpe} é ilustrativo ao demonstrar baixa utilização de telemedicina entre pessoas com deficiência, associada a escolaridade, residência urbana e barreiras de familiaridade tecnológica. Ainda que o contexto saudita não seja diretamente transponível ao brasileiro, a inferência analítica é pertinente, pois sistemas digitais de acesso à proteção social e à avaliação em saúde tendem a beneficiar mais intensamente usuários com maior capital educacional, melhor conectividade e maior capacidade de produzir documentação adequada. O mesmo sentido aparece em estudo sobre barreiras de acesso à saúde entre adultos com deficiência, no qual fatores socioeconômicos, distância, custos e suporte familiar modulam o uso efetivo dos serviços \parencite{m2024astudyofdisabilityseveritybarriersandfacilitatingfactor}. Em consequência, a promessa de desburocratização do ATESTMED pode coexistir com um viés distributivo regressivo, caso não haja desenho institucional voltado à acessibilidade, à assistência para submissão documental e à compensação de desigualdades territoriais e cognitivas.

Esse ponto é decisivo porque a perícia médica federal não é apenas uma função técnica; ela é uma capacidade estatal de triagem, reconhecimento e proteção. Quando a entrada no sistema depende da correta digitalização de laudos, da compreensão de requisitos formais e da navegação em interfaces eletrônicas, parte da carga probatória e procedimental é deslocada para o cidadão. Tal deslocamento pode ser eficiente do ponto de vista gerencial, mas não é neutro do ponto de vista republicano. \textcite{nogueira20200731t0100000100aavaliacaodadeficiencianatraje} mostra, na trajetória da avaliação da deficiência no BPC, que a evolução dos instrumentos de avaliação está intrinsecamente ligada à expansão ou restrição do acesso a direitos. A lição é aplicável ao ATESTMED, uma vez que mudanças instrumentais alteram o regime de elegibilidade material, ainda que formalmente preservem a norma jurídica. Em outras palavras, a tecnologia processual pode produzir efeitos substantivos sobre quem consegue ser reconhecido como incapaz para fins previdenciários.

A padronização, por sua vez, constitui uma ambivalência central. Em tese, padronizar critérios, campos documentais e fluxos de análise reduz heterogeneidade indevida, melhora previsibilidade e facilita controle. A literatura sobre documentação médico-administrativa e classificação em saúde sugere que instrumentos estruturados podem aumentar consistência e comparabilidade \parencite{feudtner2014pediatriccomplexchronicconditionsclassificationsyst,sa2020medicalcertificateofdisabilityadescriptiveanalysisofav}. No entanto, a padronização excessiva pode empobrecer a apreensão de casos complexos, sobretudo quando a incapacidade decorre menos de um diagnóstico isolado e mais da interação entre condição clínica, funcionalidade e exigências ocupacionais. \textcite{guo2017assessmentofaphasiaacrosstheinternationalclassification} demonstram que ferramentas digitais podem apoiar avaliações funcionais sofisticadas quando ancoradas em referenciais como a Classificação Internacional de Funcionalidade, Incapacidade e Saúde; porém isso requer desenho clínico-metodológico refinado, e não mera conversão de formulários em interface eletrônica. O risco, no ATESTMED, é confundir padronização documental com padronização substantiva da incapacidade.

Além disso, a padronização pode induzir comportamentos adaptativos por parte dos diversos atores do sistema. Médicos assistentes podem aprender quais expressões, códigos ou formatos aumentam a chance de deferimento; requerentes podem buscar documentos “mais fortes”; e a própria administração pode calibrar filtros para responder a metas de estoque ou tempo de resposta. \textcite{kelman2009performanceimprovementandperformancedysfunctionanemp} mostram, em outro setor da saúde pública, que metas de desempenho podem gerar tanto melhoria real quanto disfunções por distorção comportamental. A analogia é fecunda, pois, se o sucesso do ATESTMED for medido predominantemente por velocidade, volume processado ou redução de filas, haverá incentivo institucional para privilegiar casos mais simples, endurecer filtros automáticos ou deslocar casos complexos para canais mais lentos, produzindo uma aparência de eficiência que não coincide necessariamente com melhor decisão. A governança por indicadores, sem contrapesos qualitativos, tende a gerar otimização local e perda de validade global.

Daí a importância da auditoria. Em sistemas digitais de decisão administrativa, a auditoria não pode ser concebida apenas como revisão posterior de conformidade formal. Ela deve abranger, ao menos, quatro dimensões, que são rastreabilidade do processo decisório, consistência interdecisória, monitoramento de vieses distributivos e avaliação de acurácia substantiva. A rastreabilidade exige que cada decisão possa ser reconstruída quanto aos documentos considerados, regras aplicadas, eventuais filtros automatizados e justificativas de deferimento ou indeferimento. A consistência interdecisória demanda comparação entre casos semelhantes para identificar variabilidade não explicada. O monitoramento de vieses distributivos requer observar se determinados grupos — por território, renda, escolaridade, tipo de agravo ou perfil ocupacional — enfrentam taxas desproporcionais de indeferimento, exigência documental ou migração para perícia presencial. E a acurácia substantiva depende de mecanismos de validação por amostragem, com reavaliação independente de casos decididos documentalmente.

A literatura sobre processos auditivos e de conformidade em contextos de deficiência alerta, contudo, para o paradoxo de que práticas de auditoria desenhadas para inclusão podem, se excessivamente formalistas, produzir exclusão \parencite{deveau2011workplaceaccommodationandauditbasedevaluationprocess}. Esse alerta é particularmente pertinente à Perícia Médica Federal. Uma auditoria centrada apenas em lista de verificação documental pode reforçar a invisibilidade de situações limítrofes, nas quais a insuficiência do documento não equivale à inexistência da incapacidade. Portanto, a auditoria do ATESTMED deve combinar controle procedimental com sensibilidade clínica e social, sob pena de reproduzir, em escala digital, as limitações de modelos burocráticos tradicionais.

A governança algorítmica constitui outro eixo incontornável da discussão. Ainda que o ATESTMED não seja, em sentido estrito, um sistema plenamente automatizado de decisão, qualquer plataforma que organize filas, valide campos, aplique regras de elegibilidade, sinalize inconsistências ou priorize casos já opera uma forma de governança algorítmica. O problema não é a existência de regras computacionais, mas sua opacidade, sua rigidez e a ausência de mecanismos claros de responsabilização e prestação de contas sobre seus efeitos. Em matéria de direitos sociais, a legitimidade do uso de algoritmos depende de transparência procedimental, explicabilidade suficiente, possibilidade de contestação e supervisão humana qualificada. Isso é ainda mais importante quando se considera que a incapacidade laborativa é categoria relacional e contextual, não redutível a um escore simples ou a uma taxonomia fechada.

A literatura incluída na síntese final oferece bases indiretas para essa cautela. Estudos sobre rotas de acesso a benefícios por incapacidade mostram que trajetórias de adoecimento, trabalho e proteção social são heterogêneas e frequentemente não lineares \parencite{sainsbury2006routesontoincapacitybenefitsfindingsfromqualitati,corden2001incapacitybenefitsandworkincentives}. Se o sistema algorítmico pressupõe linearidade documental — diagnóstico, prazo, CID, assinatura, legibilidade, coerência — ele pode falhar justamente nos casos em que a vulnerabilidade é maior. O mesmo vale para condições mentais e psicossociais, cujo registro documental e reconhecimento historicamente enfrentam dificuldades de objetivação \parencite{hospital2016mentaldisabilityaretrospectivestudyofsocioclinica}. Nesses casos, a dependência de documentação padronizada pode ampliar sub-reconhecimento, sobretudo quando o cuidado assistencial é fragmentado ou quando o atestado médico não traduz adequadamente a limitação funcional.

A discussão sobre capacidade estatal, portanto, não deve ser reduzida à capacidade tecnológica. Um Estado capaz não é apenas aquele que processa mais rapidamente, mas aquele que decide com qualidade, corrige erros, aprende com evidências e distribui proteção de modo equitativo. O ATESTMED pode ampliar capacidade estatal ao liberar recursos periciais presenciais para casos mais complexos, reduzir deslocamentos desnecessários e racionalizar o fluxo de requerimentos \parencite{soares20241031t0000000000asistematicadoatestmednaconcessa}. Todavia, essa ampliação só se sustenta se houver investimento concomitante em infraestrutura de dados, interoperabilidade, formação dos peritos, desenho de protocolos clínico-documentais e canais de revisão acessíveis. A experiência da telemedicina durante e após a pandemia é mobilizada aqui apenas como evidência adjacente, ao sugerir que tecnologia sem suporte organizacional produz adoção desigual, uso defensivo e resultados heterogêneos, ponto que também deve orientar a governança do ATESTMED \parencite{ftouni2022challengesoftelemedicineduringthecovid19pandemica}.

Nesse sentido, a capacidade estatal pericial deve ser entendida como combinação de três camadas. A primeira é a capacidade cognitiva, voltada a produzir critérios claros, baseados em evidências, para definir quais casos são adequados à análise documental e quais exigem exame presencial. A segunda é a capacidade operacional, responsável por manter sistemas estáveis, seguros, acessíveis e integrados, com suporte ao usuário e ao decisor. A terceira é a capacidade reflexiva, orientada a monitorar resultados, revisar protocolos e incorporar aprendizado institucional. Sem essa terceira camada, o ATESTMED corre o risco de cristalizar erros em escala. A literatura sobre reabilitação e capacidade laborativa reforça que trajetórias funcionais são dinâmicas e que avaliações devem dialogar com possibilidades de recuperação, adaptação e retorno ao trabalho \parencite{przysada2019anassessmentofworkabilityafterthecompletionofre,deveau2011workplaceaccommodationandauditbasedevaluationprocess}. Um sistema excessivamente estático, orientado apenas à triagem inicial, pode perder a oportunidade de articular decisão previdenciária com políticas de reabilitação e acomodação laboral.

Outro aspecto relevante é a distinção entre documentação médica, elegibilidade e efetivo acesso a benefícios. \textcite{rajasulochana2026doeshavingadisabilitycertificateensurebenefit} mostram que possuir documento médico ou administrativo não garante, por si, fruição de benefícios. No caso do ATESTMED, essa dissociação é relevante porque a eficiência na análise documental não se traduz, automaticamente, em justiça material no acesso à proteção. Pode haver deferimento formalmente correto em casos mal documentados? Sim, mas também pode haver indeferimento formalmente correto em situações socialmente injustas, quando a documentação exigida não é produzível em contextos de precariedade assistencial. A decisão baseada em evidências, nesse domínio, não pode ser confundida com decisionismo documental. Evidência é sempre relação entre dado, contexto e inferência; quando o contexto é desigual, a interpretação do dado deve ser institucionalmente prudente.

A dimensão territorial também merece atenção. Estudos sobre concentração espacial da deficiência e acesso a serviços de reabilitação indicam que barreiras geográficas continuam estruturando desigualdades de cuidado \parencite{gao2019disabilityconcentrationandaccesstorehabilitationservice}. O ATESTMED pode mitigar parte dessas barreiras ao reduzir a necessidade de deslocamento para perícia presencial. Contudo, essa vantagem territorial pode ser neutralizada por desigualdades de conectividade, disponibilidade de profissionais para emissão de atestados adequados e acesso a equipamentos de digitalização. Em áreas remotas, a ausência de perícia presencial pode deixar de ser problema apenas se houver ecossistema local capaz de produzir a prova documental exigida. Caso contrário, a digitalização desloca o gargalo sem resolvê-lo.

A discussão metodológica da própria revisão reforça essa prudência interpretativa. O fato de a síntese final ter sido construída exclusivamente com 20 estudos de versão integral recuperada aumenta a confiabilidade interna da argumentação, mas também evidencia a fragmentação do campo. Há poucos estudos diretamente voltados à interface entre perícia previdenciária, saúde digital e governança decisória, e muitos dos trabalhos disponíveis abordam dimensões correlatas — telemedicina, documentação médico-administrativa de deficiência, instrumentos de avaliação, incentivos institucionais — em contextos nacionais distintos. Isso recomenda evitar transplantes normativos simplistas. Ainda assim, a convergência temática é suficiente para sustentar uma tese forte, segundo a qual a digitalização da decisão pericial só é compatível com uma agenda de decisão baseada em evidências se for acompanhada de mecanismos explícitos de validação, correção e transparência.

Em termos práticos, isso implica reconhecer que o ATESTMED deve operar como primeira camada documental de um modelo de estratificação inteligente de casos, e não como substituto universal da perícia médica. Casos de baixa complexidade documental e clínica podem ser resolvidos com segurança relativa por análise documental assíncrona; casos de média complexidade podem demandar complementação dirigida ou teleperícia síncrona; e casos de alta incerteza, especialmente aqueles envolvendo transtornos mentais, dor crônica, multimorbidade, funcionalidade flutuante ou nexo ocupacional controverso, devem permanecer sob avaliação presencial qualificada. A literatura sobre telemedicina aponta para a utilidade de modelos híbridos quando a avaliação clínica exige combinação de conveniência digital e exame direto \parencite{ftouni2022challengesoftelemedicineduringthecovid19pandemica}. No âmbito da Perícia Médica Federal, o híbrido não é sinal de insuficiência tecnológica, mas de maturidade institucional.

Por fim, a principal contribuição desta discussão é deslocar o debate do binômio simplista “celeridade versus fraude” para um quadro mais complexo de validade, erro decisório, assimetria informacional, equidade territorial e capacidade estatal. O ATESTMED deve ser compreendido como instrumento documental de estratificação decisória, e não como substituto universal da perícia médica. Sua legitimidade dependerá da capacidade institucional de combinar, de forma escalonada, análise documental assíncrona, teleperícia síncrona, perícia presencial, auditoria e revisão contínua dos critérios à luz de evidências administrativas e científicas. Dessa forma, a digitalização poderá ampliar acesso e eficiência sem reduzir a incapacidade laborativa a uma formalidade documental nem sacrificar a qualidade técnico-pericial da decisão.

\section{Proposta de redesenho do fluxo decisório do ATESTMED}
\label{sec:orgdc24455}
A proposta de redesenho apresentada nesta seção segue a lógica analítico-decisória adaptada de Secchi, partindo do diagnóstico do problema público, organizando alternativas institucionais, explicitando critérios de escolha, avaliando riscos e benefícios esperados e culminando em uma recomendação operacional para a Perícia Médica Federal \parencite{secchi2020analisedepoliticaspublicas}. Com isso, o fluxo proposto não é tratado como simples aperfeiçoamento administrativo, mas como recomendação de política pública baseada em evidências, submetida a critérios de eficiência, efetividade, equidade, suficiência probatória e viabilidade institucional.

A proposição de um redesenho do fluxo decisório do ATESTMED deve partir de um diagnóstico institucional preciso. O modelo atual ampliou capacidade de processamento, reduziu fricções de acesso e produziu efeitos relevantes sobre os grandes números da concessão por incapacidade, mas o fez sob tensão permanente entre celeridade, suficiência probatória, heterogeneidade clínica, risco de erro decisório e necessidade de preservação da legitimidade técnico-pericial \parencite{soares20241031t0000000000asistematicadoatestmednaconcessa}. Em termos de política pública, o problema não consiste em optar entre análise documental e perícia médica tradicional, como se fossem arranjos mutuamente excludentes, mas em desenhar uma arquitetura decisória escalonada, sensível ao risco, capaz de alocar o tipo de avaliação mais adequado à complexidade de cada caso. Nessa arquitetura, o ATESTMED ocupa a camada documental assíncrona; a teleperícia pode funcionar como etapa intermediária de esclarecimento; e a perícia presencial permanece indispensável quando a complexidade clínica, funcional ou probatória exigir avaliação direta. Tal arquitetura deve ser orientada por evidências, auditável, adaptativa e explicitamente comprometida com equidade de acesso, inclusive diante das barreiras de inclusão digital e das limitações inerentes às modalidades mediadas por tecnologia \parencite{ftouni2022challengesoftelemedicineduringthecovid19pandemica, albasheer2026telemedicineutilizationanddigitalinclusionamongpe}.

O primeiro princípio do redesenho consiste em substituir a lógica binária “documento suficiente versus documento insuficiente” por uma lógica de estratificação de risco decisório. Em sistemas de avaliação de incapacidade, a variabilidade entre casos não decorre apenas do diagnóstico informado, mas da interação entre natureza da condição, previsibilidade do curso clínico, consistência documental, impacto funcional, contexto ocupacional, duração alegada da incapacidade, histórico previdenciário e possibilidade de verificação indireta por registros convergentes. A literatura sobre incapacidade laboral e comportamento avaliativo de médicos peritos demonstra que a decisão não é um ato puramente mecânico, mas um processo interpretativo sujeito a variações de julgamento, enquadramento institucional e instrumentos disponíveis \parencite{steenbeek2011thedevelopmentofinstrumentstomeasuretheworkdisa, dijkstrasdatamakingupincapacityforwork}. Por isso, o fluxo do ATESTMED deve ser redesenhado para reconhecer explicitamente graus distintos de incerteza e, a partir deles, modular a intensidade da instrução.

Nesse modelo, a porta de entrada permanece digital, mas deixa de operar como filtro uniforme. O requerimento inicial deve alimentar uma etapa de triagem estruturada por risco, baseada em critérios documentais e clínico-administrativos previamente validados. Essa triagem não se confunde com decisão automatizada final; trata-se de classificação preliminar para definir a rota instrutória mais adequada. O objetivo é simples, consistindo em permitir que casos de baixa complexidade e alta confiabilidade documental sigam para deferimento documental célere, que casos de risco intermediário sejam encaminhados para complementação dirigida ou teleperícia síncrona e que casos de alto risco decisório, seja por gravidade, inconsistência, potencial litigiosidade, suspeita de inadequação documental ou necessidade de exame físico, sejam direcionados à perícia presencial.

A triagem por risco deve combinar, no mínimo, cinco dimensões. A primeira corresponde à suficiência formal do conjunto documental, envolvendo identificação do emitente, data, assinatura, legibilidade, temporalidade compatível, descrição clínica mínima, indicação do período de afastamento e coerência entre CID, narrativa clínica e limitação alegada. A segunda é a consistência longitudinal, expressa pela existência de documentos convergentes ao longo do tempo, exames complementares, prescrições, internações, reabilitação, atendimentos seriados ou outros marcadores de seguimento. A terceira é a plausibilidade funcional, de modo que, em vez de se limitar ao diagnóstico, o sistema deve capturar elementos sobre restrições de atividade e participação, em consonância com abordagens funcionais internacionalmente consolidadas \parencite{guo2017assessmentofaphasiaacrosstheinternationalclassification, nogueira20200731t0100000100aavaliacaodadeficiencianatraje}. A quarta é a complexidade clínico-social, envolvendo multimorbidade, transtornos mentais, condições crônicas complexas, dependência tecnológica, barreiras de acesso e vulnerabilidades que aumentem a probabilidade de erro por avaliação exclusivamente documental \parencite{feudtner2014pediatriccomplexchronicconditionsclassificationsyst, hospital2016mentaldisabilityaretrospectivestudyofsocioclinica, m2024astudyofdisabilityseveritybarriersandfacilitatingfactor}. A quinta é o risco sistêmico, identificado por padrões atípicos de emissão documental, recorrência de requerimentos, incongruências cadastrais, concentração geográfica ou institucional de laudos e outros sinais que recomendem auditoria reforçada.

A partir dessas dimensões, propõe-se uma matriz de decisão em camadas. Na camada inicial, o sistema realiza validações automáticas de integridade documental e consistência básica. Não se trata de substituir o juízo pericial, mas de eliminar ruído administrativo e organizar a informação para o decisor humano. Na camada seguinte, um escore de risco decisório classifica o caso em baixo, médio ou alto risco. Casos de baixo risco seriam aqueles em que há documentação formalmente íntegra, coerência clínica, previsibilidade do curso da incapacidade, limitação funcional compatível e baixa necessidade de exame físico. Nesses casos, o ATESTMED poderia produzir decisão documental favorável, com prazo de afastamento parametrizado por faixas clínicas e sujeito a auditoria amostral posterior. Casos de risco médio seriam encaminhados para exigência documental qualificada ou para teleperícia síncrona, especialmente quando a dúvida recair sobre extensão temporal, funcionalidade ou coerência narrativa, mas não houver, em princípio, necessidade incontornável de exame físico. Casos de alto risco, por sua vez, seriam remetidos diretamente à perícia presencial, com prioridade regulada pela gravidade e pelo potencial dano de uma decisão inadequada.

Esse desenho dialoga com uma evidência recorrente na literatura sobre telemedicina, segundo a qual modalidades remotas são úteis, mas não universalmente substitutivas. Revisões sistemáticas e estudos aplicados indicam que a telemedicina melhora acesso, reduz deslocamentos e pode ser efetiva em acompanhamento, educação em saúde e determinadas avaliações, porém enfrenta limites importantes quando a decisão depende de exame físico, interação clínica mais rica, avaliação contextual aprofundada ou quando há barreiras tecnológicas e comunicacionais \parencite{ftouni2022challengesoftelemedicineduringthecovid19pandemica, guo2017assessmentofaphasiaacrosstheinternationalclassification}. Essa literatura não autoriza equiparar telemedicina, teleperícia e ATESTMED, mas sustenta a prudência de um modelo híbrido e escalonado. Para o fluxo aqui proposto, isso significa que a teleperícia síncrona deve ser concebida como etapa complementar posterior à triagem documental pelo ATESTMED, destinada a esclarecer casos de incerteza intermediária, e não como componente interno da análise documental nem como sucedâneo integral da perícia presencial.

A definição dos critérios documentais merece tratamento mais rigoroso do que o usualmente observado em sistemas de concessão massificada. Um dos riscos centrais de arranjos fortemente documentais é a naturalização de que a existência de atestado equivale, por si, à demonstração suficiente de incapacidade laborativa. A literatura sobre documentação médico-administrativa, benefícios e incapacidade mostra que o documento médico é peça relevante, mas não autossuficiente; sua força probatória depende do contexto institucional, da qualidade da informação, da finalidade do documento e da articulação com critérios decisórios transparentes \parencite{sa2020medicalcertificateofdisabilityadescriptiveanalysisofav, rajasulochana2026doeshavingadisabilitycertificateensurebenefit, dijkstrasdatamakingupincapacityforwork}. Assim, propõe-se que o ATESTMED adote um protocolo nacional de suficiência documental com campos estruturados obrigatórios e taxonomia de inconsistências. Esse protocolo deve distinguir falhas sanáveis, que ensejam exigência complementar, de falhas críticas, que impõem encaminhamento pericial. Deve, ainda, incorporar descritores funcionais mínimos, evitando que a decisão se baseie apenas em rótulos diagnósticos.

A incorporação da dimensão funcional é decisiva para reduzir tanto indeferimentos injustos quanto deferimentos frágeis. Estudos sobre incapacidade e retorno ao trabalho indicam que a relação entre doença e incapacidade não é linear, sendo mediada por funcionalidade, contexto ocupacional, reabilitação, acomodações possíveis e trajetória social do segurado \parencite{gabbay2011niceguidanceonlongtermsicknessandincapacity, przysada2019anassessmentofworkabilityafterthecompletionofre, deveau2011workplaceaccommodationandauditbasedevaluationprocess}. Em consequência, o fluxo redesenhado deve prever que, sempre que possível, a documentação informe não apenas a condição clínica, mas as limitações concretas para atividades laborais relevantes. Isso não implica importar mecanicamente modelos de avaliação de deficiência para o campo da incapacidade temporária, mas sim reconhecer que decisões mais robustas exigem linguagem funcional minimamente padronizada.

Outro eixo do redesenho é a governança da inclusão digital. A literatura recente demonstra que o uso de serviços remotos em saúde é fortemente condicionado por acesso a dispositivos, conectividade, letramento digital, apoio familiar, renda, distância dos serviços e tipo de limitação funcional \parencite{albasheer2026telemedicineutilizationanddigitalinclusionamongpe, m2024astudyofdisabilityseveritybarriersandfacilitatingfactor, gao2019disabilityconcentrationandaccesstorehabilitationservice}. Em um sistema previdenciário de escala nacional, a adoção acrítica de canais digitais pode converter eficiência administrativa em exclusão silenciosa. Por isso, o fluxo decisório do ATESTMED deve conter, já na triagem, um módulo de vulnerabilidade de acesso. Segurados com baixa capacidade digital, deficiência sensorial, barreiras cognitivas, ausência de conectividade ou residência em áreas de baixa oferta tecnológica não devem ser penalizados por não conseguirem cumprir exigências digitais complexas. Nesses casos, o sistema deve oferecer rotas assistidas, apoio multicanal, possibilidade de complementação por intermédio de unidades presenciais e priorização de perícia presencial quando a mediação tecnológica comprometer a confiabilidade da instrução.

A etapa de encaminhamento para teleperícia também requer padronização metodológica. A teleperícia não pode ser mero espelhamento improvisado da consulta presencial em ambiente virtual, nem simples prolongamento do ATESTMED documental. É necessário definir requisitos técnicos, roteiro de entrevista, critérios de elegibilidade, mecanismos de autenticação, registro audiovisual quando juridicamente cabível, e parâmetros para conversão imediata em perícia presencial diante de insuficiência informacional. A literatura sobre desafios da telemedicina durante a pandemia é clara ao apontar que problemas de comunicação, perda de vínculo, limitações de exame e desigualdades tecnológicas afetam a qualidade da decisão \parencite{ftouni2022challengesoftelemedicineduringthecovid19pandemica}. Logo, a teleperícia deve ser tratada como modalidade pericial própria, com protocolo específico, treinamento dos peritos e monitoramento de desempenho comparado.

No que se refere à perícia presencial, o redesenho proposto não a reduz a instância residual, mas a reposiciona como recurso de alta resolução para casos em que o ganho informacional do encontro clínico-pericial é decisivo. Isso inclui, entre outros, quadros musculoesqueléticos com necessidade de exame físico dirigido, transtornos mentais com dúvida diagnóstica ou funcional relevante, situações de inconsistência documental, suspeita de fraude, multimorbidade complexa e casos em que a vulnerabilidade social ou comunicacional comprometa a avaliação telepericial ou presencial. A literatura qualitativa sobre trajetórias de acesso a benefícios por incapacidade mostra que processos decisórios opacos ou excessivamente padronizados tendem a produzir insegurança, percepção de arbitrariedade e afastamento entre experiência vivida e enquadramento institucional \parencite{sainsbury2006routesontoincapacitybenefitsfindingsfromqualitati, corden2001incapacitybenefitsandworkincentives}. O redesenho, portanto, deve aumentar a inteligibilidade da rota decisória para o segurado, explicando por que determinado caso foi deferido documentalmente, remetido à teleperícia ou convocado para exame presencial.

Um componente central da proposta é a auditoria em circuito fechado. Em vez de auditoria meramente repressiva ou formal, propõe-se um sistema de auditoria clínica e administrativa orientado por risco e aprendizado. Casos deferidos documentalmente devem ser submetidos a amostragem estratificada, com sobrepeso para grupos diagnósticos de maior variabilidade, emissores de alta frequência, padrões territoriais atípicos e benefícios com duração acima de determinados limiares. Casos indeferidos também devem ser auditados, pois o erro de exclusão é tão relevante quanto o erro de inclusão. A auditoria deve comparar a decisão inicial com desfechos subsequentes, como necessidade de recurso, reversão em instâncias posteriores, concessão posterior por outra via, tempo até retorno ao trabalho, reentrada no sistema e eventual judicialização. Esse desenho permite estimar sensibilidade, especificidade operacional e valor preditivo das rotas decisórias, ainda que em ambiente administrativo.

A retroalimentação por evidências constitui o elemento que transforma o fluxo redesenhado em sistema adaptativo. Os achados da auditoria devem alimentar revisões periódicas dos critérios de triagem, dos limiares de risco, dos protocolos documentais e das indicações de encaminhamento para teleperícia ou perícia presencial. Em termos de governança, isso exige um núcleo permanente de avaliação do ATESTMED, com participação da Perícia Médica Federal, áreas de dados, gestão previdenciária e, idealmente, instâncias de controle interno. O núcleo deve operar com painéis de indicadores que não se limitem à produtividade. A literatura sobre metas e desempenho em serviços públicos adverte que indicadores estreitos podem induzir disfunções, deslocando o comportamento organizacional para o cumprimento formal de números em detrimento da finalidade substantiva \parencite{kelman2009performanceimprovementandperformancedysfunctionanemp}. No contexto do ATESTMED, medir apenas tempo de decisão ou volume de requerimentos processados pode estimular deferimentos ou indeferimentos mecanizados, sem ganho real de qualidade. Por isso, o painel deve incluir indicadores indiretos de acurácia, taxa de reversão, heterogeneidade entre decisores, equidade territorial, satisfação procedimental, necessidade de reabertura e incidência de encaminhamento inadequado entre modalidades.

A proposta também demanda atenção à padronização do comportamento avaliativo. Estudos sobre médicos avaliadores mostram que diferenças individuais de interpretação, tolerância à incerteza e estilo decisório influenciam o resultado, mesmo sob regras comuns \parencite{steenbeek2011thedevelopmentofinstrumentstomeasuretheworkdisa}. Em um fluxo híbrido como o aqui proposto, a padronização não deve significar engessamento, mas redução de variabilidade indevida. Para tanto, recomenda-se a elaboração de guias decisórios por grupos de condições prevalentes, com exemplos de suficiência documental, sinais de alerta, critérios de encaminhamento e parâmetros de duração esperada, sempre preservando margem para julgamento clínico fundamentado. Tais guias devem ser continuamente recalibrados com base nos dados de auditoria e nos achados da literatura.

Há, ainda, uma dimensão territorial que não pode ser negligenciada. A distribuição desigual de serviços de saúde, reabilitação e infraestrutura digital afeta tanto a produção dos documentos quanto a capacidade do segurado de cumprir exigências e participar de teleperícias ou cumprir exigências digitais \parencite{gao2019disabilityconcentrationandaccesstorehabilitationservice, mackett2025ametareviewofliteraturereviewsofdisabilitytravel}. Um fluxo nacional uniforme, se cego a essas assimetrias, tende a reproduzir desigualdades regionais. O redesenho deve, portanto, incorporar parâmetros territoriais na triagem e no monitoramento, identificando áreas com maior taxa de exigências não cumpridas, maior conversão de teleperícia em presencial, maior indeferimento por insuficiência documental e maior tempo de resolução. Esses dados podem orientar políticas complementares de apoio, inclusive parcerias locais para facilitação de acesso.

Do ponto de vista normativo e institucional, o redesenho proposto preserva a centralidade da Perícia Médica Federal, mas redefine sua atuação em chave de inteligência regulatória. O perito deixa de ser mobilizado indistintamente para todos os casos e passa a concentrar esforço onde sua intervenção agrega maior valor informacional e maior proteção contra erro. Isso é compatível com diretrizes que enfatizam intervenção proporcional, apoio ao retorno e permanência no trabalho quando possível, e articulação entre avaliação e funcionalidade \parencite{gabbay2011niceguidanceonlongtermsicknessandincapacity}. Ao mesmo tempo, evita-se o risco de desmedicalização simplista da incapacidade, que poderia converter a análise documental em procedimento burocrático desancorado de expertise clínica.

Em síntese, o fluxo decisório redesenhado para o ATESTMED deve operar como sistema escalonado de decisão baseada em evidências. Esse sistema combinaria entrada digital universal, triagem por risco decisório, aplicação de critérios documentais estruturados, encaminhamento seletivo para deferimento documental, exigência complementar, teleperícia síncrona ou perícia presencial, auditoria clínica e administrativa orientada por risco e retroalimentação contínua dos critérios a partir de dados de desempenho e desfechos. Trata-se de uma proposta que busca conciliar eficiência, justiça procedimental, robustez probatória e adaptabilidade institucional. Seu mérito principal reside em abandonar tanto a crença de que todo caso exige presença física quanto a ilusão de que o documento, isoladamente, resolve a complexidade da incapacidade laborativa. Entre esses extremos, o que se propõe é uma arquitetura decisória proporcional à incerteza, metodologicamente auditável e permanentemente corrigida por evidências produzidas no próprio funcionamento do sistema.

\section{Indicadores, monitoramento e avaliação}
\label{sec:org46db49c}
A consolidação do ATESTMED como dispositivo de saúde digital aplicado à Perícia Médica Federal exige uma arquitetura de monitoramento e avaliação que ultrapasse a mensuração de volume processado e incorpore, de modo explícito, dimensões de qualidade decisória, equidade, tempestividade, consistência técnico-pericial, reversibilidade administrativa e judicial, auditabilidade e aprendizado institucional. Em sistemas decisórios mediados por documentos clínicos digitalizados, algoritmos de triagem, regras de negócio e fluxos documentais digitais, o desempenho não pode ser inferido apenas por produtividade ou redução de filas. Ao contrário, a literatura sobre incapacidade laboral, documentação médico-administrativa, retorno ao trabalho, telemedicina e avaliação institucional sugere, em conjunto, que intervenções organizacionais dessa natureza produzem efeitos simultaneamente desejáveis e potencialmente distorcivos, razão pela qual o desenho de indicadores deve ser multidimensional, longitudinal e sensível ao contexto \parencite{soares20241031t0000000000asistematicadoatestmednaconcessa, ftouni2022challengesoftelemedicineduringthecovid19pandemica, gabbay2011niceguidanceonlongtermsicknessandincapacity, kelman2009performanceimprovementandperformancedysfunctionanemp}.

A lógica de monitoramento proposta decorre do referencial analítico adotado no artigo, uma vez que decisões baseadas em evidências exigem rastreabilidade da informação, suficiência documental, transparência dos critérios e capacidade de revisão. Por isso, o painel de indicadores deve permitir escrutínio do fluxo decisório, identificação de vieses e retroalimentação contínua dos protocolos do ATESTMED, sem reduzir desempenho a produtividade administrativa.

Os indicadores de processo constituem a primeira camada do sistema de monitoramento. Nessa dimensão, importa medir o fluxo operacional desde o protocolo do requerimento até a conclusão administrativa, distinguindo etapas e modalidades, como recepção documental, validação cadastral, checagem de integridade do atestado, classificação de elegibilidade para análise documental, eventual encaminhamento para teleperícia síncrona ou perícia presencial, emissão da decisão e comunicação ao requerente. A utilidade desses indicadores não reside apenas em descrever o volume processado, mas em identificar gargalos, retrabalho e pontos de fricção entre tecnologia e prática pericial. Taxas de completude documental, proporção de requerimentos com anexos legíveis, percentual de documentos com metadados consistentes, frequência de exigências complementares, taxa de encaminhamento do ATESTMED para teleperícia e taxa de migração do canal documental ou telepericial para o presencial são métricas centrais, pois expressam a qualidade do insumo decisório e a adequação da modalidade escolhida. Em telemedicina e saúde digital, a literatura mostra que limitações de infraestrutura, usabilidade, padronização documental e letramento digital afetam diretamente a efetividade do serviço \parencite{ftouni2022challengesoftelemedicineduringthecovid19pandemica, albasheer2026telemedicineutilizationanddigitalinclusionamongpe}. No caso do ATESTMED, isso implica reconhecer que parte do desempenho observado pode decorrer menos da regra decisória em si e mais da heterogeneidade da documentação submetida e da capacidade do usuário de navegar o sistema.

Os indicadores de resultado, por sua vez, devem captar os efeitos imediatos e intermediários do ATESTMED sobre a concessão do benefício por incapacidade temporária e sobre a organização da Perícia Médica Federal. A taxa de deferimento, isoladamente, é insuficiente e pode induzir interpretações equivocadas. É necessário desagregá-la por perfil clínico, duração estimada do afastamento, grupo diagnóstico, território, canal de entrada, recorrência do requerente e rota instrutória adotada, seja deferimento documental, exigência complementar, teleperícia ou perícia presencial. A literatura sobre incapacidade e benefícios demonstra que trajetórias de acesso são heterogêneas e socialmente condicionadas \parencite{sainsbury2006routesontoincapacitybenefitsfindingsfromqualitati, corden2001incapacitybenefitsandworkincentives, nogueira20200731t0100000100aavaliacaodadeficiencianatraje}. Assim, um mesmo percentual agregado de concessão pode ocultar seletividades relevantes, como maior indeferimento em grupos com menor capacidade de produzir documentação clínica robusta ou em regiões com menor acesso a especialistas e serviços diagnósticos. Também devem ser acompanhados resultados organizacionais, como redução da demanda reprimida presencial, redistribuição da carga de trabalho pericial e variação do estoque de processos pendentes, em linha com os achados de impacto sistêmico já observados na implementação do ATESTMED \parencite{soares20241031t0000000000asistematicadoatestmednaconcessa}.

A qualidade decisória requer um conjunto próprio de indicadores, porque a decisão pericial não se confunde com celeridade nem com conformidade procedimental. Trata-se de aferir se a conclusão administrativa é tecnicamente consistente, suficientemente fundamentada, proporcional à evidência disponível e estável quando submetida a revisão. Estudos sobre comportamento de médicos peritos e instrumentos de mensuração da avaliação de incapacidade indicam que variabilidade interavaliador, uso desigual de critérios e diferenças na interpretação funcional são problemas estruturais em sistemas securitários \parencite{steenbeek2011thedevelopmentofinstrumentstomeasuretheworkdisa, dijkstrasdatamakingupincapacityforwork}. No ATESTMED, essa preocupação se intensifica porque a mediação documental reduz a observação clínica direta. Por isso, devem ser monitoradas a taxa de concordância entre decisão documental e decisão posterior em perícia presencial, a estabilidade da decisão em instâncias recursais, a consistência entre casos semelhantes e a aderência a protocolos clínico-funcionais previamente definidos. Sempre que possível, convém construir amostras de dupla leitura cega para estimar concordância interavaliador e detectar deriva decisória ao longo do tempo. A qualidade decisória também pode ser inferida pela suficiência da motivação registrada, mediante auditorias qualitativas de prontuários administrativos e análise textual estruturada das justificativas.

A dimensão da equidade é indispensável em qualquer avaliação séria do ATESTMED. A literatura internacional sobre deficiência, acesso a serviços e inclusão digital mostra que barreiras tecnológicas, territoriais, socioeconômicas e funcionais afetam desproporcionalmente pessoas com deficiência e grupos vulneráveis \parencite{albasheer2026telemedicineutilizationanddigitalinclusionamongpe, m2024astudyofdisabilityseveritybarriersandfacilitatingfactor, gao2019disabilityconcentrationandaccesstorehabilitationservice, mackett2025ametareviewofliteraturereviewsofdisabilitytravel}. Em consequência, o monitoramento deve comparar indicadores por sexo, faixa etária, escolaridade, marcadores indiretos de renda, localização urbana/rural, região, tipo de vínculo laboral, grupo diagnóstico e, quando disponível, marcadores de deficiência e dependência funcional. A equidade não se resume a igualdade de tratamento formal; ela exige examinar se o canal digital produz exclusão indireta. Taxas mais elevadas de exigência documental, abandono do requerimento, indeferimento por insuficiência de prova ou conversão para presencial em determinados grupos podem sinalizar barreiras de acesso e não necessariamente menor elegibilidade substantiva. A experiência de documentação e avaliação da deficiência em outros contextos reforça que a posse de documentação médica ou administrativa não garante, por si, acesso efetivo a direitos, havendo mediações institucionais que podem reproduzir exclusões \parencite{rajasulochana2026doeshavingadisabilitycertificateensurebenefit, sa2020medicalcertificateofdisabilityadescriptiveanalysisofav, hospital2016mentaldisabilityaretrospectivestudyofsocioclinica, deveau2011workplaceaccommodationandauditbasedevaluationprocess}.

Os indicadores de tempo de resposta devem ser tratados com cautela analítica. A redução do tempo entre requerimento e decisão é um objetivo legítimo, sobretudo em benefícios de natureza substitutiva de renda, mas a literatura de gestão pública adverte que metas temporais podem gerar disfunções de desempenho quando passam a orientar o comportamento organizacional de forma exclusiva \parencite{kelman2009performanceimprovementandperformancedysfunctionanemp}. No ATESTMED, o tempo mediano de decisão, o percentil 90 de conclusão, o tempo até a primeira análise e o tempo de resolução de exigências complementares são métricas relevantes, porém devem ser lidas em conjunto com indicadores de reversão, retrabalho e qualidade. Uma queda abrupta no tempo médio acompanhada de aumento de recursos, judicialização ou encaminhamentos posteriores para perícia presencial pode indicar compressão indevida da análise. Em termos de desenho avaliativo, recomenda-se o uso de painéis balanceados, nos quais indicadores de celeridade sejam sempre apresentados ao lado de métricas de acurácia e estabilidade decisória.

A reversão administrativa e a judicialização funcionam como indicadores sentinela da robustez do modelo. A taxa de reforma em recurso administrativo, a proporção de decisões documentais posteriormente alteradas em teleperícia ou perícia presencial e o percentual de concessões ou indeferimentos revertidos judicialmente devem ser acompanhados por coorte temporal e por estratos clínicos. Não se trata de assumir que toda reversão denote erro originário, pois mudanças clínicas, produção superveniente de prova e diferenças de padrão decisório entre instâncias também influenciam o desfecho. Ainda assim, padrões persistentes de reversão concentrados em certos perfis de caso podem revelar insuficiência dos critérios documentais ou inadequação da modalidade assíncrona para determinadas condições, especialmente aquelas em que a funcionalidade, a observação comportamental ou a interação clínica são decisivas. A literatura sobre capacidade para o trabalho e incapacidade de longo prazo sugere que a avaliação funcional é mais robusta quando incorpora abordagem multidimensional e, em muitos casos, componente ocupacional ou contextual \parencite{gabbay2011niceguidanceonlongtermsicknessandincapacity, przysada2019anassessmentofworkabilityafterthecompletionofre, deveau2011workplaceaccommodationandauditbasedevaluationprocess}. Isso recomenda prudência na expansão indiscriminada do modelo documental para todos os perfis de requerimento.

A auditoria deve ser concebida não como mecanismo meramente punitivo, mas como infraestrutura de validação contínua. Indicadores de auditoria incluem proporção de casos auditados, taxa de não conformidade formal, taxa de divergência técnico-decisória, reincidência de falhas por unidade ou perfil de analista, tempo de saneamento das inconformidades e impacto das recomendações corretivas. Auditorias amostrais estratificadas, combinando revisão documental, reclassificação independente e análise de padrões anômalos, são particularmente adequadas para o ATESTMED. A experiência de processos avaliativos baseados em auditoria mostra que instrumentos de conformidade podem, paradoxalmente, produzir exclusões se privilegiarem lista de verificação formal em detrimento da substância do direito material \parencite{deveau2011workplaceaccommodationandauditbasedevaluationprocess}. Por isso, a auditoria do ATESTMED deve equilibrar conformidade procedimental, suficiência probatória e justiça decisória, evitando que a busca por padronização reduza a sensibilidade a contextos clínicos complexos.

Por fim, o aprendizado institucional deve ser explicitamente mensurado, pois sistemas digitais maduros não apenas processam casos, mas aprendem com seus próprios dados. Indicadores dessa dimensão incluem frequência de revisão de protocolos, tempo entre identificação de problema e implementação de ajuste, número de ciclos de retroalimentação concluídos, produção de notas técnicas baseadas em evidências, incorporação de achados de auditoria em capacitação e evolução da concordância decisória após intervenções educacionais. A literatura sobre telemedicina e avaliação em saúde aponta, por analogia, que a melhoria sustentada depende de mecanismos reflexivos, participação dos atores e uso sistemático de evidências para redesenho de processos \parencite{ftouni2022challengesoftelemedicineduringthecovid19pandemica, steenbeek2011thedevelopmentofinstrumentstomeasuretheworkdisa}. No âmbito da Perícia Médica Federal, isso significa transformar dados operacionais em inteligência regulatória e pedagógica, com governança capaz de distinguir variação aleatória, mudança de perfil da demanda e falhas estruturais do modelo.

Em síntese, o monitoramento do ATESTMED deve operar por meio de um painel integrado de indicadores de processo, resultado, qualidade decisória, equidade, tempo de resposta, reversão, judicialização, auditoria e aprendizado institucional, articulado a uma estratégia de avaliação contínua e metodologicamente transparente. A principal exigência analítica é evitar leituras unidimensionais de sucesso. Em um sistema de decisão baseada em evidências, a efetividade do ATESTMED somente pode ser afirmada quando a celeridade administrativa vier acompanhada de consistência técnico-pericial, redução de assimetrias de acesso, estabilidade das decisões e capacidade institucional de corrigir rumos à luz de evidências produzidas com rastreabilidade documental e análise de textos integrais.

\section{Limitações}
\label{sec:org617c153}


A presente investigação apresenta limitações que devem ser explicitadas de modo rigoroso, tanto para qualificar o alcance inferencial de seus achados quanto para orientar agendas futuras de pesquisa sobre ATESTMED, saúde digital e decisão baseada em evidências na Perícia Médica Federal. A primeira limitação diz respeito ao próprio desenho da revisão e à constituição do conjunto de estudos analisados. Embora tenha sido adotado um procedimento de curadoria compatível com rastreabilidade metodológica, a síntese final foi deliberadamente restrita aos estudos cuja versão integral original foi efetivamente recuperada, em consonância com a regra previamente definida para a revisão. Nesse processo, foram lidos 248 registros no conjunto preliminar, tendo-se alcançado o alvo de 20 estudos incluídos com versão integral recuperada, exatamente o número previsto para a síntese final. Todavia, 12 estudos potencialmente elegíveis foram excluídos por versão integral não recuperada, ao passo que 216 registros deixaram de ser avaliados após o atingimento do alvo amostral. Ainda que tal decisão seja defensável do ponto de vista da padronização e da auditabilidade do processo de seleção, ela introduz uma limitação substantiva, pois a síntese não esgota o universo temático disponível, representando antes um recorte condicionado por disponibilidade documental, profundidade de busca e saturação operacional da curadoria. Em outras palavras, a robustez interna da síntese não elimina a possibilidade de viés de disponibilidade de evidência.

Essa limitação é particularmente relevante porque o campo em análise é intrinsecamente interdisciplinar e disperso entre literatura de telemedicina, avaliação de incapacidade, documentação médico-administrativa de deficiência, retorno ao trabalho, administração pública e políticas de proteção social. Em contextos dessa natureza, a recuperação do texto integral não é um detalhe técnico, mas um determinante da própria arquitetura da evidência. A exclusão de estudos sem acesso ao texto integral reduz o risco de interpretação baseada apenas em resumos, mas pode favorecer a sobrerrepresentação de periódicos, repositórios e jurisdições com maior acessibilidade digital. Ademais, a habilitação de bases como Elsevier Scopus e CORE ampliou a sensibilidade da busca, porém não elimina assimetrias de indexação, sobretudo para literatura cinzenta, documentos institucionais e estudos aplicados de circulação restrita, frequentemente relevantes em temas de gestão pública e perícia administrativa. Assim, a revisão deve ser lida como uma síntese qualificada do conjunto de estudos recuperado, e não como mapeamento exaustivo de toda a produção pertinente.

Uma segunda limitação decorre da elevada heterogeneidade dos estudos incluídos. O conjunto final reúne desenhos metodológicos, populações, contextos institucionais, objetos analíticos e desfechos muito distintos entre si. Há estudos sobre telemedicina em populações com deficiência \parencite{albasheer2026telemedicineutilizationanddigitalinclusionamongpe,ftouni2022challengesoftelemedicineduringthecovid19pandemica}, investigações sobre documentação médico-administrativa e acesso a benefícios \parencite{rajasulochana2026doeshavingadisabilitycertificateensurebenefit,sa2020medicalcertificateofdisabilityadescriptiveanalysisofav,nogueira20200731t0100000100aavaliacaodadeficiencianatraje}, análises de incapacidade laboral e retorno ao trabalho \parencite{gabbay2011niceguidanceonlongtermsicknessandincapacity,sainsbury2006routesontoincapacitybenefitsfindingsfromqualitati,corden2001incapacitybenefitsandworkincentives,przysada2019anassessmentofworkabilityafterthecompletionofre}, além de estudos sobre comportamento avaliativo de peritos e instrumentos de mensuração \parencite{steenbeek2011thedevelopmentofinstrumentstomeasuretheworkdisa,dijkstrasdatamakingupincapacityforwork}. Tal diversidade, embora intelectualmente fecunda, impõe restrições à comparabilidade direta e inviabiliza qualquer pretensão de agregação quantitativa simples. A situação é análoga àquela observada em revisões sobre afastamento prolongado e incapacidade, nas quais a heterogeneidade dos estudos limita a síntese a planos predominantemente descritivos \parencite{gabbay2011niceguidanceonlongtermsicknessandincapacity}.

No caso específico deste artigo, a heterogeneidade afeta pelo menos quatro dimensões. Primeiro, os conceitos centrais não são uniformes, pois “incapacidade”, “deficiência”, “funcionalidade”, “elegibilidade”, “telemedicina” e “inclusão digital” aparecem definidos de maneiras distintas conforme o sistema jurídico-sanitário de origem. Segundo, os desfechos variam entre acesso, satisfação, uso de serviços, retorno ao trabalho, concessão de benefícios, perfil sociodemográfico e desempenho institucional. Terceiro, os níveis de análise oscilam entre indivíduo, serviço, território e sistema de proteção social. Quarto, a qualidade metodológica dos estudos é desigual, com presença de levantamentos transversais, estudos retrospectivos, análises qualitativas, revisões e estudos instrumentais. Em consequência, qualquer inferência transversal sobre o ATESTMED deve ser formulada com prudência, evitando extrapolações causais fortes a partir de evidência estruturalmente fragmentada.

Uma terceira limitação refere-se ao uso de aproximações analíticas indiretas. Dada a escassez de estudos diretamente centrados no ATESTMED enquanto dispositivo tecnoburocrático específico da Perícia Médica Federal, parte da argumentação precisou mobilizar evidências adjacentes, como telemedicina em populações com deficiência, documentação médico-administrativa, avaliação funcional, barreiras de acesso e incentivos institucionais. Esse procedimento é metodologicamente legítimo em revisões de campos emergentes, mas implica reconhecer que nem todo achado sobre teleconsulta, documentação médico-administrativa de deficiência ou incapacidade laboral é automaticamente transferível para a lógica decisória do ATESTMED. A utilização dessas aproximações amplia a inteligibilidade do fenômeno, porém reduz a especificidade empírica da inferência. Por exemplo, barreiras de inclusão digital identificadas em usuários com deficiência \parencite{albasheer2026telemedicineutilizationanddigitalinclusionamongpe,m2024astudyofdisabilityseveritybarriersandfacilitatingfactor} são altamente plausíveis no contexto brasileiro, mas não substituem estudos observacionais sobre o perfil real dos requerentes do ATESTMED, seus padrões de uso, taxas de deferimento, necessidade de complementação documental e eventuais desigualdades regionais. De modo semelhante, discussões sobre distorções induzidas por metas e indicadores em serviços públicos \parencite{kelman2009performanceimprovementandperformancedysfunctionanemp} oferecem uma lente analítica útil para pensar efeitos não intencionais da digitalização pericial, mas permanecem como aproximações conceituais, não como demonstrações empíricas diretas sobre a Perícia Médica Federal.

A quarta limitação concerne à transferibilidade internacional da evidência. Parte expressiva dos estudos analisados provém de contextos institucionais muito distintos do brasileiro, incluindo Reino Unido, Arábia Saudita, Índia, Polônia, Portugal e outros sistemas com diferentes arranjos de seguridade social, atenção à saúde, documentação médico-pericial da incapacidade e maturidade digital \parencite{gabbay2011niceguidanceonlongtermsicknessandincapacity,albasheer2026telemedicineutilizationanddigitalinclusionamongpe,rajasulochana2026doeshavingadisabilitycertificateensurebenefit,przysada2019anassessmentofworkabilityafterthecompletionofre,sa2020medicalcertificateofdisabilityadescriptiveanalysisofav}. Essa diversidade é útil para identificar padrões gerais, como a persistência de barreiras de acesso, a importância da intervenção precoce, o papel da funcionalidade e os limites da mediação digital não presencial \parencite{ftouni2022challengesoftelemedicineduringthecovid19pandemica,gabbay2011niceguidanceonlongtermsicknessandincapacity}. Contudo, a validade externa desses achados para o caso brasileiro é necessariamente condicionada. O ATESTMED opera em um ambiente normativo, federativo e administrativo singular, marcado por desigualdades territoriais, heterogeneidade de infraestrutura digital, judicialização, assimetria informacional entre usuários e Estado e coexistência de modelos presenciais e documentais de avaliação. Portanto, a literatura internacional ilumina mecanismos e hipóteses, mas não autoriza transposição linear de resultados.

Há, ainda, uma limitação temporal e tecnológica. O campo da saúde digital evolui em ritmo acelerado, e parte dos estudos incluídos foi produzida em conjunturas específicas, especialmente durante ou logo após a pandemia de COVID-19, quando a expansão da telemedicina ocorreu sob condições excepcionais \parencite{ftouni2022challengesoftelemedicineduringthecovid19pandemica}. Isso significa que alguns achados podem refletir adaptações emergenciais, e não padrões estabilizados de uso, qualidade ou aceitabilidade. No caso do ATESTMED, cuja institucionalização se insere em trajetória própria de transformação digital do Estado, a interpretação de evidências pandêmicas deve ser feita com cautela, distinguindo efeitos conjunturais de mudanças estruturais. Além disso, tecnologias, interfaces, protocolos de autenticação e rotinas de triagem documental tendem a se modificar rapidamente, o que pode reduzir a atualidade operacional de conclusões derivadas de estudos produzidos em janelas temporais curtas.

Por fim, a limitação mais importante talvez resida nas lacunas empíricas brasileiras. Apesar da relevância crescente do tema, ainda são escassos estudos nacionais de base empírica robusta que examinem, de forma integrada, o desempenho do ATESTMED, os perfis dos requerentes, os padrões de decisão, os efeitos distributivos da digitalização, a qualidade documental dos atestados, a concordância entre avaliação documental e avaliação presencial subsequente, bem como os impactos sobre equidade, eficiência e legitimidade institucional. O estudo de \textcite{soares20241031t0000000000asistematicadoatestmednaconcessa} constitui contribuição importante, mas não esgota as múltiplas dimensões analíticas requeridas para uma avaliação abrangente do dispositivo. Também permanecem insuficientemente exploradas as interfaces entre ATESTMED e deficiência, saúde mental, condições crônicas complexas, barreiras territoriais e vulnerabilidades sociais, temas que a literatura internacional sugere serem decisivos para a qualidade da decisão em contextos de incapacidade e documentação médica \parencite{hospital2016mentaldisabilityaretrospectivestudyofsocioclinica,feudtner2014pediatriccomplexchronicconditionsclassificationsyst,gao2019disabilityconcentrationandaccesstorehabilitationservice}. Em consequência, este artigo oferece uma base analítica consistente para problematizar o tema, mas não substitui a necessidade de estudos brasileiros longitudinais, avaliativos e comparativos, com dados administrativos, desenho quase experimental quando possível, e articulação entre evidência quantitativa, análise documental e investigação qualitativa sobre práticas periciais e experiência dos usuários.

Em síntese, as limitações aqui reconhecidas não invalidam os achados, mas delimitam seu estatuto epistemológico. Trata-se de uma síntese interpretativa fundada em estudos selecionados por análise de textos integrais, metodologicamente rastreável, porém condicionado por heterogeneidade, uso de aproximações indiretas, restrições de transferibilidade internacional e insuficiência de evidência brasileira específica. Justamente por isso, a principal contribuição desta seção é reforçar que a agenda de decisão baseada em evidências na Perícia Médica Federal exige não apenas melhor uso da literatura existente, mas também produção sistemática de evidência nacional primária, sensível às singularidades institucionais, tecnológicas e sociais do ATESTMED no Brasil.

\section{Conclusão}
\label{sec:org840bdea}
A pergunta que orientou este artigo consistiu em saber em que medida o ATESTMED, compreendido como dispositivo de saúde digital aplicado à análise documental da incapacidade laborativa, pode contribuir para uma decisão baseada em evidências na Perícia Médica Federal sem comprometer a consistência técnico-pericial, a equidade de acesso e a legitimidade institucional. A resposta construída ao longo do estudo, a partir de revisão estruturada com análise estrita de textos integrais, indica que o ATESTMED representa uma inovação relevante e potencialmente estruturante, mas sua efetividade pública depende menos da mera digitalização do fluxo e mais da qualidade epistemológica, organizacional e regulatória do arranjo decisório em que se insere. Em outros termos, o ganho de escala, celeridade e racionalização administrativa observado no contexto do ATESTMED somente se converte em melhoria substantiva da decisão pericial quando articulado a critérios transparentes de elegibilidade, padronização mínima da prova documental, monitoramento de desfechos, governança de risco e mecanismos permanentes de recalibração baseados em evidências \parencite{soares20241031t0000000000asistematicadoatestmednaconcessa, steenbeek2011thedevelopmentofinstrumentstomeasuretheworkdisa}.

Do ponto de vista metodológico, a contribuição do artigo foi construir uma síntese analítico-propositiva sustentada por revisão estruturada da literatura e por análise de textos integrais. Esse desenho permitiu articular evidências sobre saúde digital, incapacidade laboral, documentação médico-administrativa, telemedicina, avaliação funcional e desempenho institucional, preservando rastreabilidade da base empírica e evitando inferências apoiadas apenas em resumos ou metadados. As evidências sobre telemedicina foram utilizadas como literatura adjacente, e não como substituto de evidências específicas sobre ATESTMED, justamente para preservar a distinção entre análise documental assíncrona, teleperícia síncrona e perícia presencial.

Os achados convergem para quatro conclusões substantivas. A primeira é que a digitalização da porta de entrada pericial, por si só, não elimina os problemas clássicos da avaliação de incapacidade; ela os reconfigura. A literatura examinada demonstra que processos de documentação médica, concessão e avaliação funcional são historicamente atravessados por heterogeneidade de critérios, assimetrias de informação, influência do contexto ocupacional e tensões entre padronização e julgamento profissional \parencite{gabbay2011niceguidanceonlongtermsicknessandincapacity, dijkstrasdatamakingupincapacityforwork, steenbeek2011thedevelopmentofinstrumentstomeasuretheworkdisa}. O ATESTMED, ao deslocar parte da decisão para a análise documental digital assíncrona, reduz fricções presenciais e amplia capacidade de processamento, mas também intensifica a dependência da qualidade do atestado, da inteligibilidade clínica do documento e da aderência entre linguagem assistencial e linguagem pericial. Assim, a inovação não substitui a necessidade de método; ao contrário, torna-a mais urgente.

A segunda conclusão é que a literatura sobre telemedicina e saúde digital oferece suporte consistente para reconhecer ganhos de acesso, continuidade e conveniência, mas igualmente adverte para limites importantes quando a decisão depende de exame físico, apreciação funcional contextualizada ou avaliação de casos complexos \parencite{ftouni2022challengesoftelemedicineduringthecovid19pandemica, guo2017assessmentofaphasiaacrosstheinternationalclassification, albasheer2026telemedicineutilizationanddigitalinclusionamongpe}. Esse ponto é central para a Perícia Médica Federal, desde que formulado com precisão, pois o ATESTMED não é teleperícia, mas análise documental digital assíncrona. Por isso, ele deve ser interpretado como primeira camada de estratificação decisória, e não como substituto universal da perícia presencial. Casos com alta suficiência documental, menor ambiguidade diagnóstica e menor necessidade de exame direto podem ser adequadamente processados pela via documental; casos com incerteza intermediária podem exigir teleperícia síncrona; e casos com maior incerteza, multimorbidade, sofrimento mental, controvérsia funcional ou repercussão ocupacional complexa exigem escalonamento para avaliação médico-pericial presencial. A racionalidade adequada, portanto, é híbrida, seletiva e orientada por risco, e não dicotômica entre “digital” e “presencial”.

A terceira conclusão é que a decisão baseada em evidências, no âmbito pericial, não pode ser reduzida à incorporação de literatura científica sobre doenças ou tecnologias. Ela requer integração entre evidência externa, dados administrativos, comportamento decisório dos peritos, características dos requerentes e efeitos institucionais dos indicadores de desempenho. Estudos sobre retorno ao trabalho, incentivos, reabilitação, documentação médica e avaliação funcional mostram que intervenções mais efetivas tendem a ser precoces, multidisciplinares e sensíveis ao contexto laboral e social \parencite{gabbay2011niceguidanceonlongtermsicknessandincapacity, przysada2019anassessmentofworkabilityafterthecompletionofre, corden2001incapacitybenefitsandworkincentives, sainsbury2006routesontoincapacitybenefitsfindingsfromqualitati}. Ao mesmo tempo, a literatura de administração pública alerta que metas quantitativas mal calibradas podem induzir disfunções de desempenho, deslocando o foco da qualidade para a mera produtividade \parencite{kelman2009performanceimprovementandperformancedysfunctionanemp}. Aplicado ao ATESTMED, isso significa que indicadores como tempo médio de análise, volume de concessões ou redução de fila são necessários, porém insuficientes. Eles devem ser acompanhados por métricas de acurácia, reversão recursal, necessidade posterior de perícia presencial, duração efetiva do benefício, perfil dos indeferimentos e distribuição territorial e sociodemográfica dos resultados.

A quarta conclusão é que a equidade digital constitui condição de legitimidade do modelo. A literatura sobre deficiência, acesso a serviços e inclusão digital evidencia que barreiras educacionais, territoriais, tecnológicas e funcionais afetam de modo desigual a capacidade de utilizar soluções remotas \parencite{albasheer2026telemedicineutilizationanddigitalinclusionamongpe, m2024astudyofdisabilityseveritybarriersandfacilitatingfactor, gao2019disabilityconcentrationandaccesstorehabilitationservice, mackett2025ametareviewofliteraturereviewsofdisabilitytravel}. Em consequência, um sistema como o ATESTMED pode ampliar acesso para muitos e, simultaneamente, produzir exclusão silenciosa para grupos com menor letramento digital, conectividade precária, limitações cognitivas, barreiras comunicacionais ou dificuldade de obtenção de documentação clínica adequada. A experiência internacional com documentação médico-administrativa e benefícios também sugere que o reconhecimento formal de uma condição não assegura, por si, acesso efetivo e justo às prestações \parencite{rajasulochana2026doeshavingadisabilitycertificateensurebenefit, sa2020medicalcertificateofdisabilityadescriptiveanalysisofav, nogueira20200731t0100000100aavaliacaodadeficiencianatraje}. Logo, a transformação digital da perícia deve ser acompanhada de desenho inclusivo, canais assistidos, teleperícia apenas quando apropriada, alternativas presenciais preservadas e monitoramento ativo de desigualdades de uso e resultado.

A incorporação da matriz analítica adaptada de Secchi reforça essa conclusão ao organizar o percurso argumentativo em diagnóstico do problema, análise de alternativas, explicitação de critérios e recomendação de redesenho institucional, conectando a revisão estruturada da literatura à lógica prática da decisão pública baseada em evidências \parencite{secchi2020analisedepoliticaspublicas}.

A contribuição específica deste artigo reside, portanto, em propor uma leitura do ATESTMED não apenas como ferramenta de desburocratização, mas como objeto de governança clínica e decisória. Ao reunir evidências de saúde digital, avaliação de incapacidade, documentação médico-administrativa e políticas de retorno ao trabalho, o estudo demonstra que a modernização da Perícia Médica Federal precisa ser concebida como arquitetura de decisão baseada em evidências, e não como simples automação procedimental. Essa distinção é decisiva. Automação sem governança tende a reproduzir opacidades em maior escala; governança baseada em evidências permite transformar escala em qualidade institucional. Nessa perspectiva, o ATESTMED deve operar com protocolos dinâmicos de suficiência documental, taxonomias de risco, auditoria amostral, retroalimentação dos peritos, interoperabilidade com dados assistenciais e avaliação periódica de impacto. Também se mostra recomendável o desenvolvimento de instrumentos padronizados para observar variabilidade decisória e comportamento avaliativo, à semelhança do que a literatura internacional já vem propondo para médicos avaliadores \parencite{steenbeek2011thedevelopmentofinstrumentstomeasuretheworkdisa}.

Em síntese final, retomando a pergunta inicial, os achados indicam que o ATESTMED pode, sim, fortalecer a decisão baseada em evidências na Perícia Médica Federal, desde que seja institucionalmente tratado como mecanismo de triagem qualificada e decisão proporcional à suficiência da prova, e não como solução universal para toda demanda por reconhecimento de incapacidade. O artigo contribui ao demonstrar, com base em estudos selecionados por análise de textos integrais, que os benefícios de celeridade e escala são reais, mas dependem de desenho regulatório, monitoramento analítico e compromisso com equidade e qualidade pericial. A recomendação central para a Perícia Médica Federal é consolidar o ATESTMED como modelo híbrido, estratificado por risco e permanentemente avaliado por indicadores de acesso, consistência, acurácia e justiça distributiva. Somente assim a inovação digital deixará de ser apenas expediente de gestão de demanda para se converter em política pública madura de decisão pericial orientada por evidências, capaz de combinar eficiência administrativa, segurança técnica e proteção social legítima \parencite{soares20241031t0000000000asistematicadoatestmednaconcessa, ftouni2022challengesoftelemedicineduringthecovid19pandemica, gabbay2011niceguidanceonlongtermsicknessandincapacity}.

\section{Referências}
\label{sec:org4191d69}
\printbibliography[heading=none]
\end{document}